\pdfoutput=1
\documentclass[12pt]{article}

\usepackage{amsmath}
\usepackage{amssymb}
\usepackage{amsthm}
\usepackage{mathtools}
\usepackage[margin=2.5cm]{geometry}
\usepackage[utf8]{inputenc}
\usepackage{lmodern}
\usepackage{microtype}
\usepackage{url}
\usepackage{hyperref}

\theoremstyle{plain}
\newtheorem{theorem}{Theorem}[section]
\newtheorem{lemma}[theorem]{Lemma}
\newtheorem{proposition}[theorem]{Proposition}
\newtheorem{corollary}[theorem]{Corollary}
\newtheorem{conjecture}[theorem]{Conjecture}

\theoremstyle{definition}
\newtheorem{definition}[theorem]{Definition}

\theoremstyle{remark}
\newtheorem{remark}[theorem]{Remark}

\newcommand{\ord}{\operatorname{ord}}
\newcommand{\lcm}{\operatorname{lcm}}
\newcommand{\Nm}{\operatorname{N}}

\title{Chebyshev primality criteria for Wagstaff numbers}
\author{Alexey Dolotov\thanks{Email: \texttt{science@dolotov.com}}}
\date{}

\begin{document}

\maketitle

\begin{abstract}
The Chebyshev base \(\omega_3 = 3 + 2\sqrt{2}\) and the identity \(W_p + 1 = 4W_{p-2}\) yield an unconditional \(N+1\) primality criterion for the Wagstaff numbers \(W_p = (2^p + 1)/3\); combined with a companion \(N-1\) criterion and the Brillhart--Lehmer--Selfridge framework, this gives fully verified primality proofs of \(W_{2617}\), \(W_{10501}\), and \(W_{12391}\), independent of ECPP. For the sufficiency conjecture --- whether the single congruence \(\omega_3^{(W_p+1)/2} \equiv -1 \pmod{W_p}\) implies primality --- we show that Condition~(I) holds automatically, exclude factors \(r \equiv 5, 7 \pmod{8}\), reduce the split branch \(r \equiv 1 \pmod{8}\) to the nonexistence of compatible Pell triples (proved for rank parameters \(d \leq 81\)), and prove that the inert branch \(r \equiv 3 \pmod{8}\) is impossible for order parameter \(d \leq 43\). An Order-Pinning Lemma shows that every inert prime factor \(r\) of a composite \(W_p\) satisfies \(\operatorname{ord}_r(2) = 2p\), and a Multi-Factor Pinning Theorem forces all inert factors of a common \(W_p\) to share a single order parameter; combined with a multi-factor bound, these unconditionally exclude the entire finite catalog of small-\(d\) danger configurations produced by the parity obstruction. An exhaustive Platinum Lemma rules out every inert factor with \(r \leq 10^{12}\) via a \(684\,965\,381\)-row enumeration, and an elementary deterministic Second-Moment Reduction transfers the residual problem to the single claim that no composite \(W_p\) admits a pair of distinct residual inert danger factors \(r_1, r_2 > 10^{12}\) at the same~\(p\).

\end{abstract}

\bigskip\noindent
\textit{2020 Mathematics Subject Classification.} Primary 11Y11; Secondary 11A51, 11B39.

\smallskip\noindent
\textit{Key words and phrases.} Wagstaff numbers, primality testing, Chebyshev polynomials, Pell sequences, Brillhart--Lehmer--Selfridge.

\section{Introduction}


The Wagstaff numbers \(W_p = (2^p + 1)/3\) for odd primes \(p\) are the
natural ``$+1$'' companions of the Mersenne numbers \(M_p = 2^p-1\).
For Mersenne numbers, the Lucas--Lehmer test \cite{ref12} gives a proven
standalone criterion (see \cite{ref17} for the historical development of
Lucas-type tests and \cite{ref7} for a modern computational treatment). For
Wagstaff numbers, by contrast, published primality proofs have
relied on general-purpose methods such as elliptic curve primality proving (ECPP), while the known
Chebyshev/Lucas-type tests remained conjectural on the sufficiency
side.

The paper delivers two distinct results, summarised in \S\S1.1--1.2
below. Sections~2--3 establish the unconditional criteria and proofs~(A);
Sections~4--7 develop the case-by-case reduction for~(B); Section~8
adds the cross-case unconditional reductions from pinning; Section~9
presents the computational verification.

\textbf{(A) Three unconditional primality proofs.} We prove
\(W_{2617}\), \(W_{10501}\), and \(W_{12391}\) prime by
Chebyshev/Brillhart--Lehmer--Selfridge (BLS) methods, independent of ECPP (Theorems~\ref{thm:W2617}--\ref{thm:W12391}).

\textbf{(B) Reduction of the standalone sufficiency conjecture
to a single quantitative residual.} The sufficiency of the single
Chebyshev congruence \(\omega_3^{(W_p+1)/2} \equiv -1 \pmod{W_p}\)
(Conjecture~\ref{conj:sufficiency}) is reduced to two explicit unconditional lower
bounds on multiplicative orders --- one at split factors, one at
inert factors --- with every known Wagstaff prime and every tested
composite \(W_p\) consistent with the conjecture (\S\S9.1, 9.5). The
inert-factor branch is further reduced by an Order-Pinning Lemma,
Multi-Factor Pinning Theorem, Platinum Lemma, and a deterministic
Second-Moment Reduction to the single claim that no composite
\(W_p\) admits a pair of distinct residual inert danger factors
\(r_1, r_2 > 10^{12}\) at the same~\(p\) (\S8).

Contribution (A) does not depend on Conjecture~\ref{conj:sufficiency} and may be
read separately from (B).

\subsection{Three new unconditional primality proofs}


The base \(a = 3\) in the Chua framework gives a universal \(N+1\)
condition for Wagstaff primes through \(\omega_3 = 3 + 2\sqrt{2}\).
The identity \(W_p + 1 = 4W_{p-2}\) turns this into an unconditional
primality criterion (Theorem~\ref{thm:criterion}), and a companion \(N-1\) condition
via \(a = 5\) yields a combined \(N+1/N-1\) criterion when
\(p \equiv 1 \pmod{6}\) (Theorem~\ref{thm:combined}). Combining these criteria
with the Brillhart--Lehmer--Selfridge framework \cite{ref3}, we
give fully verified primality proofs of

\begin{itemize}
\item \(W_{2617}\) (788 digits, Theorem~\ref{thm:W2617}),
\item \(W_{10501}\) (3161 digits, Theorem~\ref{thm:W10501}), and
\item \(W_{12391}\) (3730 digits, Theorem~\ref{thm:W12391}),
\end{itemize}

independent of the ECPP proofs in the literature. These three
results are entirely unconditional and do not rely on the
Chebyshev sufficiency conjecture (Conjecture~\ref{conj:sufficiency}) discussed
below.

\subsection{Structural reduction of Conjecture~\ref{conj:sufficiency}}\label{sec:status-table}


The second contribution concerns the sufficiency of Condition (II)
on a hypothetical composite \(W_p\) (Conjecture~\ref{conj:sufficiency}) --- the
question of whether the single Chebyshev congruence
\(\omega_3^{(W_p+1)/2} \equiv -1 \pmod{W_p}\) suffices, on its
own, to decide primality of \(W_p\). The conjecture is open,
but we reduce it to a small number of explicit local problems
and verify those problems for every known Wagstaff prime and
every analysed factor configuration.

Condition (I) in the Chebyshev test is automatically satisfied at every
prime factor of \(W_p\), whether \(W_p\) is prime or composite
(Theorem~\ref{thm:cond-I}). The standalone problem therefore reduces entirely to
Condition (II).

Prime factors \(r \equiv 5, 7 \pmod{8}\) cannot divide \(W_p\)
(Lemma~\ref{lem:mod-8-exclusion}). A local \(2\)-adic order classification for
\(\omega_3\) (Lemma~\ref{lem:v2-classification}) classifies
\(v_2(\operatorname{ord}_r(\omega_3))\) in every residue class
\(r \bmod 8\) --- exactly at inert factors and by a sharp upper
bound at split factors --- and supplies the exact local input
used in \S6 to isolate the inert branch.

The remaining hypothetical composite counterexamples are analysed
branch by branch, by the residue class of the prime factors
\(r \mid W_p\) modulo 8:

\begin{itemize}
\item \(r \equiv 5, 7 \pmod{8}\): \textbf{impossible} (Lemma~\ref{lem:mod-8-exclusion}).
\item \(r \equiv 1 \pmod{8}\): \textbf{reduced to Conjecture~\ref{conj:R3}} via Theorem~\ref{thm:R3-conditional}.
  Conjecture~\ref{conj:R3} is a Pell-rank divisibility statement; via
  Proposition~\ref{prop:rank-equivalence}, Conjecture~\ref{conj:R3} is equivalent to the nonexistence of
  compatible triples (NCT), which is proved for all rank
  parameters \(d \leq 81\) (Theorem~\ref{thm:NCT-81}).
  Conjecture~\ref{conj:R3} is verified computationally at all 185
  known split factors (\S9.2).
\item \(r \equiv 3 \pmod{8}\): \textbf{order parameter \(d \leq 43\) proved
  impossible} (Theorem~\ref{thm:pell-exclusion}). Among 869 structurally analysed
  inert factors, 845 of 869 satisfy \(G_r \le 43\) and are closed
  unconditionally by Corollary~\ref{cor:small-Gr}. An exhaustive inert-factor
  survey over 684{,}965{,}381 pairs \((p,r)\) with
  \(r < 10^{12}\) and \(p < 5 \times 10^{11}\) finds no composite
  counterexample. A natural subclass of the surviving
  configurations (those inert primes whose quotient \((r+1)/4\)
  has every odd prime factor in the set \(A\) of Class III primes;
  see \S6.2) has density zero within the primes
  \(r \equiv 3 \pmod{8}\) (Proposition~\ref{prop:density-zero}); a naive Hooley-type
  GRH argument does \emph{not} extend this to a pointwise closure
  (\S10, Problem 2).
\end{itemize}

\textbf{Cross-case unconditional reductions from pinning.} Beyond the
case-by-case analysis above, Section~\ref{sec:pinning-reductions}
establishes a sequence of unconditional reductions that apply
uniformly to every candidate composite \(W_p\):

\begin{itemize}
\item Every inert prime factor \(r \mid W_p\) with \(r > 3\) satisfies
  \(\operatorname{ord}_r(2) = 2p\) (Order-Pinning Lemma,
  Theorem~\ref{thm:order-pinning}).
\item All inert prime factors of a single composite \(W_p\) share a
  common order parameter \(p\) (Multi-Factor Pinning,
  Theorem~\ref{thm:multi-factor-pinning}).
\item Consequently, the complete finite catalog of danger
  configurations with \(d \leq 1000\) produced by the parity
  classification of \S\ref{sec:nct} is \emph{unconditionally}
  excluded (Phantom Exclusion,
  Theorem~\ref{thm:phantom-exclusion}): within a single composite, no
  two inert factors can both have small \(d\).
\item An exhaustive Platinum enumeration removes every inert factor
  with \(r \leq 10^{12}\) (Theorem~\ref{thm:platinum}).
\item A deterministic second-moment reduction
  (Theorem~\ref{thm:second-moment-reduction}) transfers the
  residual all-inert branch to the single claim that no composite
  \(W_p\) admits two distinct inert prime factors \(r_1, r_2\)
  with \(r_i > 10^{12}\) and \(d_{r_i} > 43\) --- the sole
  remaining unconditional input for the all-inert branch of
  Conjecture~\ref{conj:sufficiency}.
\end{itemize}

Appendix A records a short package of unconditional local results
at split Wagstaff factors (a Pell-form quartic character of 2,
a Wagstaff-specific Octic Dichotomy, a genus formula, and a
closed-form QNR Dichotomy in terms of \(r \bmod 16\)). These are
independent of the main reduction of \S5 and may be read
separately.

The status of each component, in summary:

\begin{center}
\small
\begin{tabular}{@{}p{5.5cm}p{5.6cm}l@{}}
\hline
Component & Status & Reference \\
\hline
BLS proof of \(W_{2617}\), \(W_{10501}\), \(W_{12391}\) & Unconditional & Thm~\ref{thm:W2617}--\ref{thm:W12391} \\
Cond.\ (I) at every \(r \mid W_p\) & Unconditional & Thm~\ref{thm:cond-I} \\
\(r \equiv 5, 7\) excluded & Unconditional & Lem~\ref{lem:mod-8-exclusion} \\
\(r \equiv 1\): Conj.~\ref{conj:R3} \(\Rightarrow\) failure & Unconditional & Thm~\ref{thm:R3-conditional} \\
\(r \equiv 1\), NCT for \(d \leq 81\) & Unconditional & Thm~\ref{thm:NCT-81} \\
\(r \equiv 1\), NCT for \(83 \leq d \leq 200\) & Computational$^*$ & \S9.4 \\
\(r \equiv 1\), Conj.~\ref{conj:R3} (all 185 known) & Computational$^*$ & \S9.2 \\
\(r \equiv 3\), \(G_r \leq 43\) & Unconditional & Cor~\ref{cor:small-Gr} \\
Refined SFP mod 16 & Unconditional & Thm~\ref{thm:sfp-mod16} \\
Order-Pinning: \(\operatorname{ord}_r(2) = 2p\) for \(r \mid W_p\), \(r > 3\) & Unconditional & Thm~\ref{thm:order-pinning} \\
Multi-Factor Pinning (all factors share \(p\)) & Unconditional & Thm~\ref{thm:multi-factor-pinning} \\
Phantom Exclusion of full catalog for \(d \leq 1000\) & Unconditional & Thm~\ref{thm:phantom-exclusion} \\
Multi-factor bound: \(\geq\!3\) inert factors per \(W_p\) & Unconditional & Thm~\ref{thm:multi-factor} \\
Platinum Lemma: no danger factor with \(r \leq 10^{12}\) & Computational$^*$ (\(684\,965\,381\) rows) & Thm~\ref{thm:platinum} \\
Second-Moment Reduction & Unconditional & Thm~\ref{thm:second-moment-reduction} \\
\(r \equiv 3\), residual tail pair-count \((r > 10^{12})\) & Open & \S10, Prob.~2 \\
\(r \equiv 3\), \(G_r > 43\) & Natural subclass density zero (uncond.); pointwise open & Prop~\ref{prop:density-zero} \\
Full Conj.~\ref{conj:sufficiency} & Open & \S10 \\
\hline
\multicolumn{3}{@{}l}{\small $^*$Exhaustive finite search, deterministic and reproducible from the supplementary scripts.}\\
\end{tabular}
\end{center}

All scripts, factorization files, BLS certificates, and verification
logs referenced in this paper are archived on Zenodo
(\href{https://doi.org/10.5281/zenodo.19496206}{10.5281/zenodo.19496206});
see \S9.7 for the software toolchain and ``Supplementary material''
at the end of the References for the full inventory.

\subsection{What remains open}


The full sufficiency of Condition (II) (Conjecture~\ref{conj:sufficiency}) remains open.
The unconditional gap consists of two structurally independent problems:

\begin{enumerate}
\item \textbf{Conjecture~\ref{conj:R3} at split factors:} at every split prime divisor
  \(r \equiv 1 \pmod 8\), Conjecture~\ref{conj:R3} says
  \(p \mid \operatorname{ord}_r(1+\sqrt 2)\). It is known to hold
  in three explicit cases (Cases A, B, and \(4U_p^2 \equiv 1\); see
  \S5.3), and in general it is equivalent (Proposition~\ref{prop:rank-equivalence}) to the
  nonexistence of compatible Pell triples (NCT), which we prove for
  all \(d \leq 81\) (Theorem~\ref{thm:NCT-81}) and verify computationally
  for the 11 parity-unblocked values \(d \in [83, 200]\) (\S9.4).
  All 185 known split factors satisfy Conjecture~\ref{conj:R3}: 29\%
  by Cases A/B/\(4U\), 23\% by rank-parity exclusion
  (Proposition~\ref{prop:rank-equivalence}), and the remaining 48\% by
  direct computation.
  The asymptotic gap (large \(d\)) is open.
\item \textbf{Tail pair-count for inert factors.} After the
  case-by-case reduction of \S6 and the cross-case pinning
  reductions of \S\ref{sec:pinning-reductions}, the inert branch is
  closed except for a single quantitative residual. By
  Order-Pinning and Multi-Factor Pinning the inert factors of any
  composite \(W_p\) all share a common order parameter \(p\); by
  the Phantom Exclusion Theorem the finite catalog of danger
  configurations with \(d \leq 1000\) produced by the parity
  classification is unconditionally excluded; by the Platinum
  Lemma no inert factor with \(r \leq 10^{12}\) can participate in
  a composite \(W_p\); and by the Second-Moment Reduction
  (Theorem~\ref{thm:second-moment-reduction}) the number of
  all-inert composite counterexamples is at most the total
  unordered pair count of residual inert danger factors
  \((r_1, r_2)\) of a single \(W_p\) with
  \(r_i > 10^{12}\), \(\operatorname{ord}_{r_i}(2) = 2p\), and
  \(d_{r_i} > 43\).

  \textbf{Residual Problem (tail pair count).} Show that this
  total pair count vanishes: no composite \(W_p\) admits two
  distinct inert prime factors meeting these conditions. This
  vanishing, together with the split-branch problem above, implies
  Conjecture~\ref{conj:sufficiency} in full.
\end{enumerate}

Both problems lie in the regime of unconditional lower bounds on
multiplicative orders, where Artin-type questions are well known
to be hard. The contribution of (B) is that the gap is
structurally narrow: after the pinning reductions of
\S\ref{sec:pinning-reductions}, all \(684\,965\,381\) inert
factors of composite \(W_p\) tested in \S9.5 are consistent with
the conjecture, together with all 36 known Wagstaff primes and
probable primes (\S9.1).

\subsection{Notation}


Throughout, \(p \geq 5\) denotes an odd prime and
\(W_p = (2^p+1)/3\) the corresponding Wagstaff number; \(r\) denotes
a prime divisor of \(W_p\). For an integer or algebraic integer
\(x\) coprime to a modulus \(m\), \(\operatorname{ord}_m(x)\) denotes
the multiplicative order of \(x\) modulo \(m\), and \(v_q(\cdot)\)
denotes the \(q\)-adic valuation. We work in
\(\mathbb{Z}[\sqrt{2}]\), with the fundamental unit
\(\alpha = 1 + \sqrt{2}\) and the Chebyshev base
\(\omega_3 = 3 + 2\sqrt{2} = \alpha^2\). The integer Pell sequences
\((U_n)_{n \geq 0}\) and \((V_n)_{n \geq 0}\) are defined by
\(V_n = \alpha^n + \bar{\alpha}^n\) (where
\(\bar{\alpha} = 1 - \sqrt{2}\)) and
\(\alpha^n = V_n/2 + U_n\sqrt{2}\); equivalently they satisfy the
common recurrence \(X_{n+1} = 2X_n + X_{n-1}\) with initial values
\(U_0 = 0\), \(U_1 = 1\) and \(V_0 = V_1 = 2\), and
\(V_n^2 - 8U_n^2 = 4(-1)^n\). For a prime \(r\)
with \(r \neq 2\) and \((2/r) = 1\), \(\rho(r)\) denotes the rank
of apparition of \(r\) in \((U_n)\), i.e. the least \(n \geq 1\)
with \(r \mid U_n\). For an inert prime factor \(r\) of \(W_p\)
we write \(G_r := \gcd((r+1)/4,\, W_{p-2})\), and the ``order
parameter'' of \(r\) is
\[
  d_r \;=\; g \;:=\; \operatorname{ord}_r(\omega_3)/4
  \;=\; (r+1)/(4\,\gcd((r+1)/4,\,W_{p-2}));
\]
we use \(d_r\) when emphasising dependence on the factor and the
bare symbol \(g\) inside local proofs where no ambiguity arises;
both refer to the same integer. See \S6.1 for the role these
quantities play in the gap structure.
\(\Phi_d(x)\) is the \(d\)-th cyclotomic polynomial, and \(A\)
denotes the set of odd primes \(q\) with
\(v_2(\operatorname{ord}_q(2)) = 1\); following the classification
introduced in \S6.2, these are the ``Class~III'' primes.



\section{Chebyshev bases for Wagstaff numbers}


\subsection{The Chua framework}


For an integer \(a \neq 0, \pm 1\), define \(\omega_a = a + \sqrt{a^2-1}\),
so \(\omega_a \bar{\omega}_a = 1\). The identity
\(\omega_a^n + \omega_a^{-n} = 2T_n(a)\), where \(T_n\) is the
\(n\)-th Chebyshev polynomial of the first kind, motivates the
name ``Chebyshev base.''

\begin{theorem}[{Chua \cite{ref5}; cf.\ Williams \cite[Ch.~4]{ref17}}]\label{thm:chua}
Let \(Q\) be an odd prime with
\(\gcd(a^2-1, Q) = 1\). Set \(D = a^2-1\), \(\varepsilon = (D/Q)\),
\(\delta = (2(a+1)/Q)\). Then
\(\omega_a^{(Q-\varepsilon)/2} \equiv \delta \pmod{Q}\)
in \(\mathbb{Z}[\sqrt{D}]/(Q)\).
\end{theorem}


\begin{proof}
Define \(\gamma = a + 1 + \sqrt{D}\). Then
\(\gamma^2 = 2(a+1) \cdot \omega_a\). Compute \(\gamma^Q\) in
\(\mathbb{Z}[\sqrt{D}]/(Q)\) by Frobenius
(\(\gamma^Q = a + 1 + \varepsilon\sqrt{D}\)) and via the square
(\(\gamma^Q = \delta \cdot \omega_a^{(Q-1)/2} \cdot \gamma\)).
Equating and dividing by \(\gamma\) (a unit) yields
\(\omega_a^{(Q-\varepsilon)/2} = \delta\).\end{proof}


\subsection{Basic properties of Wagstaff numbers}


\begin{lemma}[Order of 2]\label{lem:ord-2}
For every prime \(p \geq 5\) and every
prime \(r \mid W_p\): \(\operatorname{ord}_r(2) = 2p\) and
\(r \equiv 1 \pmod{2p}\).
\end{lemma}


\begin{proof}
Since \(3W_p = 2^p + 1\): \(2^p \equiv -1 \pmod{r}\).
Therefore \(\operatorname{ord}_r(2) \mid 2p\) but
\(\operatorname{ord}_r(2) \nmid p\). Since \(p\) is prime,
\(\operatorname{ord}_r(2) \in \{2, 2p\}\). If \(\operatorname{ord} = 2\):
\(r \mid 3\), contradicting \(\gcd(W_p, 3) = 1\)
(since \(p \geq 5\) is prime, \(p \nmid 6 = \operatorname{ord}_9(2)\),
so \(v_3(2^p+1) = 1\) and \(3 \nmid W_p\)). So
\(\operatorname{ord}_r(2) = 2p\), and \(2p \mid r-1\) by Fermat.\end{proof}


\begin{lemma}\label{lem:mod-8-exclusion}
For all odd primes \(p \geq 3\):
\textup{(i)}~\(W_p \equiv 3 \pmod{8}\);
\textup{(ii)}~every prime factor of \(W_p\) satisfies
\(r \equiv 1\) or \(3 \pmod{8}\);
\textup{(iii)}~the total multiplicity of prime factors
\(\equiv 3\) is odd (i.e.\ \(\sum_{r \mid W_p,\; r \equiv 3\,(8)} v_r(W_p)\) is odd).
\end{lemma}

\begin{proof}
(i)~Since \(p\) is odd, \(3W_p = 2^p + 1 \equiv 1 \pmod{8}\),
so \(W_p \equiv 3 \pmod{8}\).
(ii)~By Lemma~\ref{lem:ord-2}, \(\operatorname{ord}_r(2) = 2p\)
with \(v_2(2p) = 1\). For \(r \equiv 7\): \((2/r) = 1\), so
\(\operatorname{ord}_r(2)\) divides the odd number \((r-1)/2\),
contradicting \(v_2 = 1\). For \(r \equiv 5\):
\(v_2(\operatorname{ord}_r(2)) = v_2(r-1) = 2 \neq 1\).
(iii)~Write \(W_p = \prod r_i^{e_i}\). By (ii), each
\(r_i \equiv 1\) or \(3 \pmod{8}\). Since
\(1^e \equiv 1\) and \(3^e \equiv 3^e \pmod{8}\), the product
is \(\equiv 3^{\sum' e_i} \pmod{8}\) where \(\sum'\) runs over
\(r_i \equiv 3\). From \(W_p \equiv 3 \pmod{8}\): \(\sum' e_i\)
is odd.\end{proof}

\begin{theorem}[Refined split-factor parity modulo 16]\label{thm:sfp-mod16}
For every odd prime \(p \geq 5\):
\begin{itemize}
\item[\textup{(a)}] \(W_p \equiv 11 \pmod{16}\).
\item[\textup{(b)}] Every prime factor \(r > 3\) of \(W_p\) with
  \(r \equiv 3 \pmod 8\) satisfies \(r \equiv 3\) or
  \(r \equiv 11 \pmod{16}\).
\item[\textup{(c)}] Suppose that \(W_p\) is composite and passes
  Condition~(II) at every prime factor; then all factors are inert
  (by Theorem~\ref{thm:R3-conditional}) and the total count of
  inert factors counted with multiplicity, \(t\), and the subtotal
  of those \(\equiv 11 \pmod{16}\), \(k_{11}\), satisfy
  \[
    t + 2\,k_{11} \;\equiv\; 3 \pmod{4}.
  \]
\end{itemize}
\end{theorem}

\begin{proof}
(a) For \(p \geq 5\): \(2^5 = 32\), and the multiplicative order of
\(2\) in \((\mathbb{Z}/96\mathbb{Z})^\times\) equals \(2\), so
\(2^p \equiv 32 \pmod{96}\) for every prime \(p \geq 5\). Thus
\(2^p + 1 \equiv 33 \pmod{96}\), and dividing by \(3\) (an invertible
residue modulo~\(32\)) gives \(W_p \equiv 11 \pmod{32}\), hence
\((a)\) modulo \(16\).

(b) By Lemma~\ref{lem:mod-8-exclusion}(ii), every prime divisor
\(r > 3\) of \(W_p\) satisfies \(r \equiv 1\) or \(3 \pmod{8}\);
reducing modulo \(16\) splits \(r \equiv 3 \pmod 8\) into
\(r \equiv 3\) or \(r \equiv 11 \pmod{16}\).

(c) Under the stated hypotheses every factor is inert
\((\equiv 3 \pmod{8})\). Write the product with multiplicities,
\(W_p = \prod r_i\). By (b), each \(r_i \in \{3, 11\}
\pmod{16}\). In \((\mathbb{Z}/16\mathbb{Z})^\times\) the element
\(3\) has order \(4\) with
\(\{3^0, 3^1, 3^2, 3^3\} = \{1, 3, 9, 11\}\), so \(11 \equiv 3^3
\pmod{16}\); the product satisfies
\[
  \prod r_i \;\equiv\; 3^{\,k_3 + 3\,k_{11}} \pmod{16}
\]
where \(k_3\) (resp.\ \(k_{11}\)) counts factors \(\equiv 3\)
(resp.\ \(\equiv 11\)) modulo~\(16\), and
\(k_3 + k_{11} = t\). Using (a) and \(11 \equiv 3^3 \pmod{16}\),
\(k_3 + 3\,k_{11} \equiv 3 \pmod 4\); substituting \(k_3 = t -
k_{11}\) gives \(t + 2\,k_{11} \equiv 3 \pmod 4\).\end{proof}

\begin{corollary}[Mod-16 sub-configurations for \(t=3\)]\label{cor:sfp-mod16-t3}
In the situation of Theorem~\ref{thm:sfp-mod16}(c), if \(t = 3\)
then \(k_{11} \in \{0, 2\}\): either all three inert factors are
\(\equiv 3 \pmod{16}\), or exactly two are \(\equiv 11\) and one is
\(\equiv 3 \pmod{16}\).
\end{corollary}
\begin{proof}
Immediate from \(3 + 2 k_{11} \equiv 3 \pmod 4\), i.e.\
\(k_{11}\) even, combined with \(0 \leq k_{11} \leq 3\).\end{proof}

\begin{lemma}[{$v_2$ classification for $\omega_3$}]\label{lem:v2-classification}
Let \(r \geq 5\)
be a prime. Then
\[
v_2(\operatorname{ord}_r(\omega_3)) =
\begin{cases}
2 & \text{if } r \equiv 3 \pmod{8}, \\
1 & \text{if } r \equiv 5 \pmod{8}, \\
0 & \text{if } r \equiv 7 \pmod{8}, \\
\leq v_2(r-1)-1 & \text{if } r \equiv 1 \pmod{8}.
\end{cases}
\]
\end{lemma}


\begin{proof}
Set \(\alpha = 1+\sqrt{2}\), so \(\omega_3 = \alpha^2\) and
\(v_2(\operatorname{ord}(\omega_3)) = \max(v_2(\operatorname{ord}(\alpha))-1,\,0)\).

\emph{Split cases} (\(r \equiv 1, 7 \pmod{8}\)): \((2/r) = 1\), so
\(\sqrt{2} \in \mathbb{F}_r\) and \(\alpha \in \mathbb{F}_r^*\) with
\(\operatorname{ord}(\alpha) \mid r-1\), hence
\(v_2(\operatorname{ord}(\alpha)) \leq v_2(r-1)\).
\begin{itemize}
\item \(r \equiv 7\): \(v_2(r-1) = 1\), so \(v_2(\operatorname{ord}(\alpha)) \leq 1\);
  either way \(v_2(\operatorname{ord}(\omega_3)) = 0\).
\item \(r \equiv 1\): \(v_2(r-1) \geq 3\), giving
  \(v_2(\operatorname{ord}(\omega_3)) \leq v_2(r-1)-1\).
\end{itemize}

\emph{Inert cases} (\(r \equiv 3, 5 \pmod{8}\)): \((2/r) = -1\), so
\(\sqrt{2} \notin \mathbb{F}_r\) and \(\alpha \in \mathbb{F}_{r^2}^*\).
The Frobenius \(x \mapsto x^r\) sends \(\sqrt{2} \mapsto -\sqrt{2}\),
so \(\alpha^r = 1-\sqrt{2} = \bar{\alpha}\) and
\(\alpha^{r+1} = \alpha\bar{\alpha} = 1 - 2 = -1\). Hence
\(\operatorname{ord}(\alpha) \mid 2(r+1)\) but
\(\operatorname{ord}(\alpha) \nmid (r+1)\), which forces
\(v_2(\operatorname{ord}(\alpha)) = v_2(r+1) + 1\).
\begin{itemize}
\item \(r \equiv 3\): \(v_2(r+1) = 2\), so
  \(v_2(\operatorname{ord}(\alpha)) = 3\) and
  \(v_2(\operatorname{ord}(\omega_3)) = 2\).
\item \(r \equiv 5\): \(v_2(r+1) = 1\), so
  \(v_2(\operatorname{ord}(\alpha)) = 2\) and
  \(v_2(\operatorname{ord}(\omega_3)) = 1\).
\end{itemize}
\end{proof}

\begin{remark}
The exact value at \(r \equiv 3\) (Frobenius forces
\(\alpha^{r+1} = -1\)) is used in \S6.1 to establish
\(\operatorname{ord}(\omega_3) = 4g\) with \(g\) odd at
inert factors; the bound at \(r \equiv 1\) is used in \S5.2.
\end{remark}

\subsection{The base a = 3}


\begin{proposition}\label{prop:base-3}
For every Wagstaff prime \(W_p\) with \(p \geq 5\),
\(\omega_3 = 3 + 2\sqrt{2}\) satisfies
\(\omega_3^{(W_p+1)/2} \equiv -1 \pmod{W_p}\).
\end{proposition}


\begin{proof}
For \(a = 3\) we have \(D = a^2 - 1 = 8\) and
\(2(a+1) = 8\), so the two Legendre symbols appearing in
Theorem~\ref{thm:chua} happen to coincide:
\(\varepsilon = (D/W_p) = (8/W_p)\) and
\(\delta = (2(a+1)/W_p) = (8/W_p)\). Since
\(W_p \equiv 3 \pmod{8}\), the second supplement to quadratic
reciprocity gives
\(\left(\tfrac{2}{W_p}\right) = (-1)^{(W_p^2-1)/8}\), and
\(W_p^2 \equiv 9 \pmod{16}\) makes \((W_p^2 - 1)/8\) odd, so
\((2/W_p) = -1\). Hence
\(\varepsilon = \delta = (8/W_p) = (2/W_p)^3 = -1\). By
Theorem~\ref{thm:chua}: \(\omega_3^{(W_p+1)/2} \equiv -1\).\end{proof}


\begin{remark}[Uniqueness]\label{rem:uniqueness}
The base \(a = 3\) is the unique integer \(a \geq 2\) with
\(a^2 - 1\) a power of~2 (since \((a-1)(a+1) = 2^k\) forces
\(\{a-1,a+1\} = \{2,4\}\); cf.\ \cite{ref14}), so \(a = 3\) gives
the only universal \(N+1\) condition of this form.
\end{remark}


\subsection{The base a = 2 and the p mod 6 dichotomy}


\begin{proposition}\label{prop:base-2}
For \(\omega = 2 + \sqrt{3}\) and odd prime
\(p \geq 5\): if \(p \equiv 5 \pmod{6}\), then
\(\omega^{(W_p-1)/2} \equiv -1\) (\(N-1\) side); if
\(p \equiv 1 \pmod{6}\), then \(\omega^{(W_p+1)/2} \equiv 1\)
(\(N+1\) side).
\end{proposition}


\begin{proof}
For \(a = 2\): \(D = 3\), \(2(a+1) = 6\). By quadratic
reciprocity: for \(p \equiv 5 \pmod{6}\),
\(W_p \equiv 11 \pmod{12}\), so \((3/W_p) = +1\) and
\((6/W_p) = -1\); for \(p \equiv 1 \pmod{6}\),
\(W_p \equiv 7 \pmod{12}\), so \((3/W_p) = -1\) and
\((6/W_p) = +1\). By Theorem~\ref{thm:chua} with \(a = 2\):
\(\omega^{(W_p - \varepsilon)/2} = \delta\).\end{proof}


\subsection{The base a = 5}


\begin{proposition}\label{prop:base-5}
For \(a = 5\), \(\omega_5 = 5 + 2\sqrt{6}\),
and prime \(p \equiv 1 \pmod{6}\) with \(W_p\) prime:
\(\omega_5^{(W_p-1)/2} \equiv -1 \pmod{W_p}\).
\end{proposition}


\begin{proof}
For \(a = 5\): \(D = 24\), \(2(a+1) = 12\). For
\(p \equiv 1 \pmod{6}\): \((24/W_p) = (8/W_p)(3/W_p) = (-1)(-1) = +1\)
and \((12/W_p) = (4/W_p)(3/W_p) = (1)(-1) = -1\). By Theorem~\ref{thm:chua}:
\(\omega_5^{(W_p - 1)/2} = -1\).\end{proof}




\section{Factorization-based primality criteria}


\subsection{Order lemma}


\begin{lemma}\label{lem:order}
Let \(N\) be odd with \(N \equiv 3 \pmod{8}\). Suppose
\(\omega_3^{N+1} \equiv 1\) and \(\omega_3^{(N+1)/2} \equiv -1 \pmod{N}\).
Then for every prime \(r \mid N\): \(v_2(\operatorname{ord}_r(\omega_3))
= v_2(N+1) = 2\).
\end{lemma}


\begin{proof}
The congruences descend from \(\bmod N\) to \(\bmod r\).
Let \(d = \operatorname{ord}_r(\omega_3)\). From (I):
\(\omega_3^{N+1} \equiv 1 \pmod{r}\), so \(d \mid N+1\). From (II):
\(\omega_3^{(N+1)/2} \equiv -1 \not\equiv 1 \pmod{r}\) (since
\(r\) is an odd prime), so \(d \nmid (N+1)/2\). Together:
\(v_2(d) = v_2(N+1)\). Since \(N \equiv 3 \pmod{8}\):
\(N + 1 \equiv 4 \pmod{8}\), giving \(v_2(N+1) = 2\).\end{proof}


\subsection{General primality criterion}


\begin{theorem}\label{thm:criterion}
Let \(p \geq 5\) be prime, \(N = W_p\). Suppose
Conditions (I) \(\omega_3^{N+1} \equiv 1\) and
(II) \(\omega_3^{(N+1)/2} \equiv -1 \pmod{N}\) hold. Let \(S\) be a
set of odd prime divisors of \(W_{p-2}\) with
\(\omega_3^{(N+1)/q} - 1\) a unit mod \(N\) for each \(q \in S\).
Define \(F = 4 \prod_{q \in S} q^{v_q(W_{p-2})}\). If \(F > \sqrt{N}\),
then \(N\) is prime.
\end{theorem}


\begin{proof}
Suppose \(N\) is composite with smallest prime factor
\(r \leq \sqrt{N}\).

\textbf{Step 1 (Order constraints).} By Lemma~\ref{lem:order},
\(d := \operatorname{ord}_r(\omega_3)\) satisfies \(d \mid N+1\) and
\(v_2(d) = v_2(N+1) = 2\) (since \(N+1 = 4W_{p-2}\) with \(W_{p-2}\)
odd).

\textbf{Step 2 (Full multiplicity: \(F \mid d\)).} For each \(q \in S\),
the unit condition \(\omega_3^{(N+1)/q} - 1 \in
(\mathbb{Z}[\sqrt{2}]/(N))^{\times}\) implies, after reduction mod
\(r\), that \(\omega_3^{(N+1)/q} \not\equiv 1 \pmod{r}\). Since
\(d \mid N+1\) but \(d \nmid (N+1)/q\):
\(v_q(d) = v_q(N+1) = v_q(W_{p-2})\).
Together with \(v_2(d) = 2 = v_2(4)\):
\(F = 4 \prod q^{v_q(W_{p-2})}\) divides \(d\).

\textbf{Step 3 (Group-theoretic bound).} Since
\(\operatorname{Norm}(\omega_3) = 1\)
and \(r\) is an odd prime:
\begin{itemize}
\item If \((2/r) = -1\): \(\mathbb{Z}[\sqrt{2}]/(r) \cong \mathbb{F}_{r^2}\),
  and \(\omega_3\) lies in \(\mu_{r+1}\), so \(d \mid r+1\).
\item If \((2/r) = +1\): \(\mathbb{Z}[\sqrt{2}]/(r) \cong
  \mathbb{F}_r \times \mathbb{F}_r\), and \(d \mid r-1\).
\end{itemize}

\textbf{Case 1} (\(d \mid r-1\)): \(F \mid d \mid r-1\), so
\(r \geq F + 1 > \sqrt{N} + 1\). Contradicts \(r \leq \sqrt{N}\).

\textbf{Case 2} (\(d \mid r+1\)) --- \textbf{boundary analysis.} From
\(F \mid d \mid r+1\): \(r \geq F - 1\). Since
\(N \equiv 3 \pmod{8}\) (Lemma~\ref{lem:mod-8-exclusion}), \(N\) is not a perfect square, so
\(\sqrt{N}\) is irrational. With \(F > \sqrt{N}\) and \(F\) an integer:
\(F \geq \lfloor\sqrt{N}\rfloor + 1\). Combined with
\(r \leq \lfloor\sqrt{N}\rfloor\):
\[
\lfloor\sqrt{N}\rfloor \geq r \geq F - 1 \geq \lfloor\sqrt{N}\rfloor,
\]
so \(r = \lfloor\sqrt{N}\rfloor\) and \(d = r + 1\).

Write \(N = r \cdot s\) with \(s \geq r\). From
\((r+1) \mid (rs + 1)\) and \(r \equiv -1 \pmod{r+1}\):
\(s \equiv 1 \pmod{r+1}\).
Since \(s \geq r\), the smallest candidate is \(s = r + 2\);
the next, \(s = 2r+3\), gives
\(N = r(2r+3) = 2r^2 + 3r > r^2 + 2r + 1 = (r+1)^2\),
contradicting \(r = \lfloor\sqrt{N}\rfloor\). So
\[
N = r(r+2), \qquad N + 1 = (r+1)^2.
\]
Since \(N + 1 = 4W_{p-2}\): \(W_{p-2} = ((r+1)/2)^2\). But
\(W_{p-2} \equiv 3 \pmod{8}\) (Lemma~\ref{lem:mod-8-exclusion}), and a perfect square is
\(\equiv 0, 1,\) or \(4 \pmod{8}\). Contradiction.
\end{proof}

\begin{remark}\label{rem:bls-general}
Theorem~\ref{thm:criterion} uses the threshold \(F > \sqrt{N}\), which
suffices for small \(p\) where \(W_{p-2}\) is fully factored
(\S9.1). For larger \(p\), the factored portion satisfies
\(F > N^{1/3}\) but not \(F > N^{1/2}\). The original BLS theorem
\cite{ref3} covers this case: if \(F > N^{1/3}\), then \(N\) has at
most two prime factors, and the two-factor case is excluded by
verifying that a certain discriminant is not a perfect square.
The three primality proofs below (Theorems~\ref{thm:W2617}--\ref{thm:W12391})
all invoke this general BLS criterion, not Theorem~\ref{thm:criterion}.
\end{remark}

\subsection{Coprimality}


\begin{corollary}\label{cor:coprimality}
Every prime divisor of \(W_{p-2}\) is coprime
to \(W_p\).
\end{corollary}


\begin{proof}
\(3W_p - 12W_{p-2} = (2^p+1) - 4(2^{p-2}+1) = -3\), so
\(W_p - 4W_{p-2} = -1\).\end{proof}


\subsection{The case \texorpdfstring{\(W_{p-2}\)}{W\_\{p-2\}} prime}


\begin{corollary}\label{cor:Wp2-prime}
If \(W_{p-2}\) is prime, then
\(W_p\) is prime iff \(\omega_3^{(W_p+1)/2} \equiv -1 \pmod{W_p}\).
\end{corollary}


\begin{proof}
Necessity: Proposition~\ref{prop:base-3}. Sufficiency: let \(r \mid N\).
By Lemma~\ref{lem:order}, \(d := \operatorname{ord}_r(\omega_3)\) has
\(v_2(d) = 2\) and \(d \mid N + 1 = 4W_{p-2}\). The divisors of
\(4W_{p-2}\) with \(v_2 = 2\) are \(4\) and \(4W_{p-2}\). We exclude
\(d = 4\): \(\omega_3^4 = 577 + 408\sqrt{2}\), and
\(r \mid \gcd(576, 408) = 24\), contradicting
\(\gcd(W_p, 24) = 1\). So \(d = N + 1\). Since \(d \mid r \pm 1\):
\(r \geq N\), forcing \(r = N\).\end{proof}


\subsection{Combined N+1/N-1 criterion}


\begin{theorem}[Combined criterion]\label{thm:combined}
Let \(p \equiv 1 \pmod{6}\),
\(N = W_p\). Suppose \(\omega_3^{(N+1)/2} \equiv -1\) and
\(\omega_5^{(N-1)/2} \equiv -1 \pmod{N}\). Let \(S_+\) (resp. \(S_-\))
be sets of odd prime divisors of \(W_{p-2}\) (resp. \((N-1)/2\)) with
unit conditions analogous to Theorem~\ref{thm:criterion}. Define
\(F_+ = 4 \prod_{q \in S_+} q^{v_q(W_{p-2})}\),
\(F_- = 2 \prod_{q \in S_-} q^{v_q((N-1)/2)}\).
If \(F_+ \cdot F_- > 2N\), then \(N\) is prime.
\end{theorem}


\begin{proof}
For composite \(N\) with smallest prime factor
\(r \leq \sqrt{N}\): \(F_+ \mid \operatorname{ord}_r(\omega_3)
\mid r^2 - 1\) and \(F_- \mid \operatorname{ord}_r(\omega_5) \mid
r^2 - 1\). (For the second divisibility: \(\omega_5 \in
\mathbb{Z}[\sqrt{6}]\) has \(\operatorname{Norm}(\omega_5) = 1\),
so \(\operatorname{ord}_r(\omega_5)\) divides \(r - (6/r)\) and
hence \(r^2 - 1\).) Since \(\gcd(F_+, F_-) = 2\) (as
\(F_+ \mid N+1\) and \(F_- \mid N-1\)):
\(\operatorname{lcm}(F_+, F_-) = F_+ F_- / 2 > N > r^2 - 1\).
Contradiction.
\end{proof}


\subsection{BLS proof of \texorpdfstring{\(W_{2617}\)}{W\_2617}}


The three primality proofs in \S\S3.6--3.8 invoke the
Brillhart--Lehmer--Selfridge (BLS) framework with specific
factorizations of \(N-1\) drawn partly from the Cunningham project
tables and partly from direct computation. The exact software
toolchain --- SymPy \texttt{factorint} (with Pollard \(\rho\), \(p\!-\!1\)
and ECM substeps), gmpy2 for modular arithmetic, and PARI/GP's
APR-CL implementation \cite{ref6} for primality of
large cofactors --- and the version numbers used to produce the
certificates \texttt{bls\_certificate\_w*.json} are documented in \S9.7.
The conjecture analysis of \S\S5--6 does not depend on these proofs.

\begin{center}
\small
\begin{tabular}{lccccc}
\hline
 & Digits & \(p-1\) divisors & Primes in \(F\) & \(F > N^{1/3}\) margin & Time \\
\hline
\(W_{2617}\) & 788 & 16 & 22 & 15 bits & \(<1\) s \\
\(W_{10501}\) & 3161 & 48 & 103 & 1086 bits & 387 s \\
\(W_{12391}\) & 3730 & 32 & 61 & 953 bits & 67 s \\
\hline
\end{tabular}
\end{center}

\begin{theorem}\label{thm:W2617}
The Wagstaff number \(W_{2617}\) is prime.
\end{theorem}

\begin{proof}
Set \(N = W_{2617}\), a 788-digit number. We apply the BLS
\(N-1\) criterion \cite{ref3} (Remark~\ref{rem:bls-general}).
The exponent \(p - 1 = 2616 = 2^3 \cdot 3 \cdot 109\) has 16
divisors; the cyclotomic decomposition yields 12 factorable terms
\(\Phi_d(2)\) with \(d \leq 654\), contributing 22 prime factors
whose product \(F\) exceeds \(N^{1/3}\) by 15 bits. BLS witnesses
\(a \in \{2, 3\}\) and discriminant check complete the proof.
Condition~(II) is independently verified:
\(\omega_3^{(N+1)/2} \equiv -1 \pmod{N}\); full data in
\texttt{bls\_certificate\_w2617.json}.
\end{proof}

\subsection{BLS proof of \texorpdfstring{\(W_{10501}\)}{W\_10501}}

\begin{theorem}\label{thm:W10501}
The Wagstaff number \(W_{10501}\) is prime.
\end{theorem}

\begin{proof}
Of the 48 cyclotomic factors \(\Phi_d(2)\), 43 are factored
(38 by direct computation, 5 from the Cunningham tables
\cite{ref4}), contributing 103 primes verified by APR-CL
\cite{ref6}. BLS witnesses \(a \in \{2, 3, 5\}\) and
discriminant check complete the proof; full data in
\texttt{bls\_certificate\_w10501.json}.
\end{proof}

\subsection{BLS proof of \texorpdfstring{\(W_{12391}\)}{W\_12391}}

\begin{theorem}\label{thm:W12391}
The Wagstaff number \(W_{12391}\) is prime.
\end{theorem}

\begin{proof}
Of the 32 cyclotomic factors, 27 are factored (19 by direct
computation, 7 from the Cunningham tables \cite{ref4} and
public factor databases, 1 itself prime), contributing 61 primes
(the largest a 371-digit cofactor of \(\Phi_{2065}(2)\)), each
independently verified by APR-CL \cite{ref6}. BLS witnesses \(a \in \{2, 3\}\) and discriminant
check complete the proof; full data in
\texttt{bls\_certificate\_w12391.json}.
\end{proof}

\begin{remark}[Limits of the method]\label{rem:bls-ceiling}
The BLS \(N-1\) approach
requires \(p - 1\) to be sufficiently smooth so that enough
cyclotomic factors \(\Phi_d(2)\) have known factorizations.
Among all 36 known Wagstaff primes and probable primes
(proved for \(p \leq 42737\); see \S9.1) with \(p > 2617\), the exponents \(p = 10501\) and \(p = 12391\) are
the only ones with sufficiently smooth \(p - 1\). The next candidate
by smoothness, \(p = 14479\) (\(p - 1 = 2 \cdot 3 \cdot 19 \cdot 127\)),
falls 3956 bits short of the BLS threshold. No other known Wagstaff
exponent beyond 12391 is BLS-feasible.
\end{remark}



\section{Condition (I) and the sufficiency conjecture}


\begin{theorem}\label{thm:cond-I}
For every odd prime \(p \geq 5\) and every prime
\(r \mid W_p\):
\[
2^{(W_p-1)/2} \equiv -1 \pmod{r}.
\]
\end{theorem}


\begin{proof}
By Lemma~\ref{lem:ord-2}, \(\operatorname{ord}_r(2) = 2p\). We show
\((W_p - 1)/(2p)\) is an odd integer. Indeed:
\((W_p - 1)/2 = (2^{p-1} - 1)/3\), and
\((W_p - 1)/(2p) = (2^{p-1} - 1)/(3p)\).
This is an integer (\(p \mid 2^{p-1}-1\) by Fermat, \(3 \mid 2^{p-1}-1\)
since \(p\) is odd) and is odd (the denominator \(3p\) is odd;
the numerator \(2^{p-1}-1\) is odd since \(p \geq 5\) gives
\(2^{p-1} \equiv 0 \pmod{2}\)).
Therefore \(2^{(W_p-1)/2} = (2^p)^{(W_p-1)/(2p)} \equiv (-1)^{\text{odd}}
= -1\).\end{proof}


\begin{corollary}[CRT consequence]\label{cor:cond-I-crt}
If \(W_p\) is squarefree, then
\(2^{(W_p-1)/2} \equiv -1 \pmod{W_p}\), i.e., Condition (I) holds
for \(W_p\) itself.
\end{corollary}


\begin{proof}
Theorem~\ref{thm:cond-I} gives the congruence modulo every prime
\(r \mid W_p\). CRT yields the congruence modulo \(W_p\).\end{proof}


\begin{remark}
Squarefreeness of \(W_p\) is not known for general \(p\). All
subsequent sections use only the per-prime version
(Theorem~\ref{thm:cond-I}), so the results of \S\S5--7 are
unconditional regardless of whether \(W_p\) is squarefree.
\end{remark}

We say a composite \(W_p\) ``passes Condition~(II)'' if
\(\omega_3^{(W_p+1)/2} \equiv -1 \pmod{W_p}\). For this to hold,
it must hold modulo every prime \(r \mid W_p\).

\begin{conjecture}[Full sufficiency]\label{conj:sufficiency}
No composite \(W_p\) passes
Condition (II).
\end{conjecture}


\subsection{The two branches}


We organise the analysis by the residue class of each prime
factor \(r \mid W_p\) modulo 8. By Lemma~\ref{lem:mod-8-exclusion},
every prime factor satisfies \(r \equiv 1\) or \(3 \pmod{8}\), so
the analysis reduces to two branches:

\begin{center}
\begin{tabular}{lll}
\hline
Branch & Factors & Status \\
\hline
\textbf{Split} & \(r \equiv 1 \pmod{8}\) & Reduced to NCT (\S5) \\
\textbf{Inert} & \(r \equiv 3 \pmod{8}\) & \(d \leq 43\) proved; pure Class~III subclass density zero (\S6) \\
\hline
\end{tabular}
\end{center}




\section{The \texorpdfstring{\(r \equiv 1 \pmod{8}\)}{r == 1 (mod 8)} branch}


For \(r \equiv 1 \pmod{8}\) with \(r \mid W_p\), the algebra
\(\mathbb{Z}[\sqrt{2}]/(r)\) splits and the descent of Condition (II)
yields an odd parameter \(d\) with \(\operatorname{ord}(\omega_3) = 4d\)
and \(d \mid W_{p-2}\). After the descent of \S5.2 and the explicit
Cases A, B, and \(4U_p^2 \equiv 1\) of \S5.3, \S\S5.4--5.6 identify
\(d\) with the Pell rank \(\rho(r)/4\) (Proposition~\ref{prop:rank-equivalence}), convert
Conjecture~\ref{conj:R3} to the nonexistence of compatible triples (NCT),
and prove NCT for \(d \leq 81\) (Theorem~\ref{thm:NCT-81}).
Conjecture~\ref{conj:R3} is verified at all 185 known split factors (\S9.2).

\subsection{Notation}


We use the notation of \S1.4 throughout.

\subsection{Descent in the split algebra}


For \(r \equiv 1 \pmod{8}\) with \(r \mid W_p\): \((2/r) = 1\), so
\(\mathbb{Z}[\sqrt{2}]/(r) \cong \mathbb{F}_r \times \mathbb{F}_r\)
via \(\sqrt{2} \mapsto (s, -s)\) with \(s^2 = 2\). The element
\(\omega_3 = 3+2\sqrt{2}\) maps to \((\alpha_0, \alpha_0^{-1})\)
where \(\alpha_0 = 3+2s\).

\textbf{Descent of Condition (II).} In the split algebra,
\(\omega_3^{(N+1)/2} = -1\) becomes \(\alpha_0^{(N+1)/2} = -1\)
in \(\mathbb{F}_r^*\). Therefore:
\begin{itemize}
\item \(\operatorname{ord}(\alpha_0) \mid N+1 = 4W_{p-2}\);
\item \(v_2(\operatorname{ord}(\alpha_0)) = v_2(N+1) = 2\).
\end{itemize}

Write \(\operatorname{ord}(\alpha_0) = 4d\) with \(d\) odd and
\(d \mid W_{p-2}\).

\begin{theorem}[{Order of $\sqrt{2}$}]\label{thm:ord-sqrt2}
Let \(r \equiv 1 \pmod{8}\)
be a prime with \(r \mid W_p\). Then
\(\operatorname{ord}_{\mathbb{F}_r^*}(s) = 4p\).
\end{theorem}


\begin{proof}
\(s^{2p} = 2^p \equiv -1 \pmod{r}\), so
\(\operatorname{ord}(s) \nmid 2p\); and \(s^{4p} = 1\) gives
\(\operatorname{ord}(s) \mid 4p\). Since \(p\) is prime, the
divisors of \(4p\) that do not divide \(2p\) are exactly \(4\) and
\(4p\). The case \(\operatorname{ord}(s) = 4\) would force
\(s^4 \equiv 1 \pmod r\), i.e.\ \(4 \equiv 1 \pmod r\), contradicting
\(r \geq 5\). Hence \(\operatorname{ord}(s) = 4p\).\end{proof}


\begin{remark}
Since \(\alpha_0 = (1+s)^2\),
\(\operatorname{ord}(\alpha_0) =
\operatorname{ord}(1+s)/\gcd(\operatorname{ord}(1+s), 2)\).
The divisibility \(p \mid \operatorname{ord}(1+s)\) is the content
of Conjecture~\ref{conj:R3}.
\end{remark}

\begin{conjecture}[Pell-rank divisibility]\label{conj:R3}
For every prime \(r \equiv 1 \pmod{8}\) with
\(r \mid W_p\):
\[
p \mid \operatorname{ord}_{\mathbb{F}_r^*}(1+s),
\quad\text{where } s = \sqrt{2} \bmod r.
\]
Verified for all 185 known such primes.
\end{conjecture}


\begin{lemma}\label{lem:p-nmid-Wp2}
\(p \nmid W_{p-2}\) for every prime \(p \geq 5\).
\end{lemma}


\begin{proof}
By Fermat, \(2^p \equiv 2 \pmod{p}\), so
\(4W_{p-2} = W_p + 1 = (2^p+1)/3 + 1 = (2^p+4)/3 \equiv 2 \pmod{p}\).
Since \(p \geq 5\): \(p \nmid 2\), hence \(p \nmid W_{p-2}\).\end{proof}


\begin{theorem}[Conditional on Conjecture~\ref{conj:R3}]\label{thm:R3-conditional}
Let \(W_p\) be composite with a prime factor
\(r \equiv 1 \pmod{8}\). If Conjecture~\ref{conj:R3} holds at \(r\),
then \(r\) cannot satisfy Condition~(II).
\end{theorem}


\begin{proof}
By Conjecture~\ref{conj:R3}: \(p \mid \operatorname{ord}(1+s)\), hence
\(p \mid \operatorname{ord}(\alpha_0)\) (since \(\alpha_0 = (1+s)^2\)).
From the descent: \(\operatorname{ord}(\alpha_0) = 4d\) with
\(d \mid W_{p-2}\). So \(p \mid 4d\), hence \(p \mid d \mid W_{p-2}\),
contradicting Lemma~\ref{lem:p-nmid-Wp2}.\end{proof}


\subsection{Proved cases of Conjecture~\ref{conj:R3}}


\begin{proposition}[Algebraic identity]\label{prop:algebraic-identity}
In \(\mathbb{Z}[\sqrt{2}]\):
\[
\Phi_{2p}(\omega_3) = U_p \cdot (1+\sqrt{2})^{p-1}.
\]
Consequently \(N(\Phi_{2p}(\omega_3)) = U_p^2\).
\end{proposition}


\begin{proof}
Since \(\Phi_{2p}(x) = (x^p + 1)/(x+1)\) for prime \(p\):
\[
\omega_3^p + 1 = U_p(4U_p + V_p\sqrt{2}).
\]
The claim \((4U_p + V_p\sqrt{2})/(4+2\sqrt{2}) = \alpha^{p-1}\)
follows by induction: write \(\alpha = 1+\sqrt{2}\) and
\(f_n = (4U_n + V_n\sqrt{2})/(4+2\sqrt{2})\). The
Pell recurrences \(U_{n+1} = 2U_n + U_{n-1}\),
\(V_{n+1} = 2V_n + V_{n-1}\) give
\(4U_{n+1} + V_{n+1}\sqrt{2}
= 2(4U_n + V_n\sqrt{2}) + (4U_{n-1}+V_{n-1}\sqrt{2})\),
so \(f_{n+1} = 2 f_n + f_{n-1}\). Since
\(\alpha^2 = 2\alpha + 1\), the powers \(\alpha^{n-1}\) satisfy
the same recurrence; the initial values
\(f_0 = (\sqrt{2}-1) = \alpha^{-1}\) and
\(f_1 = 1 = \alpha^0\) match, completing the induction.\end{proof}


\begin{theorem}[Case A]\label{thm:case-A}
Let \(r \equiv 1 \pmod{8}\) with
\(r \mid W_p\) and \(r \mid U_p\). Then
\(\operatorname{ord}(\alpha) = 4p\) in the split algebra.
\end{theorem}


\begin{proof}
Reducing Proposition~\ref{prop:algebraic-identity} modulo \(r\):
\(\Phi_{2p}(\omega_3) \equiv 0\), so \(\omega_3 \bmod r\) is a
primitive \(2p\)-th root of unity, giving
\(\operatorname{ord}(\omega_3) = 2p\). Since \(\omega_3 = \alpha^2\)
and \(2p\) is even:
\(\operatorname{ord}(\alpha) = 2 \cdot \operatorname{ord}(\omega_3)
= 4p\).\end{proof}


\begin{theorem}[Case B]\label{thm:case-B}
Let \(r \equiv 1 \pmod{8}\) with
\(r \mid W_p\) and \(r \mid V_p\). Then
\(\operatorname{ord}(\alpha) = 2p\) in the split algebra.
\end{theorem}


\begin{proof}
With \(\sqrt 2 \mapsto (s, -s)\) and \(r \mid V_p\), the
identity \(\alpha^p = V_p/2 + U_p\sqrt{2}\) gives
\(\alpha^p \equiv (U_p s,\; -U_p s)\) in
\(\mathbb{F}_r \times \mathbb{F}_r\). Squaring each coordinate:
\(\alpha^{2p} \equiv (2U_p^2,\; 2U_p^2)\). From
\(V_p^2 - 8U_p^2 = -4\) with \(r \mid V_p\): \(8U_p^2 \equiv 4\),
hence \(2U_p^2 \equiv 1\), and \(\alpha^{2p} \equiv (1, 1) = 1\).
Meanwhile \(\alpha^p \equiv (U_p s,\; -U_p s)\) has its two
coordinates opposite in sign, so \(\alpha^p \neq (1, 1)\). Hence
\(\operatorname{ord}(\alpha) \mid 2p\) and \(\operatorname{ord}
(\alpha) \nmid p\), so \(\operatorname{ord}(\alpha) \in \{2, 2p\}\).
The case \(\operatorname{ord}(\alpha) = 2\) would force
\(\omega_3 = \alpha^2 \equiv (1, 1)\), i.e.\ \(3 + 2s \equiv 1\)
and \(3 - 2s \equiv 1\) in \(\mathbb F_r\); adding the two gives
\(6 \equiv 2\), i.e.\ \(4 \equiv 0 \pmod r\), contradicting
\(r \geq 17\). Therefore \(\operatorname{ord}(\alpha) = 2p\).\end{proof}


\begin{theorem}\label{thm:case-4U}
Let \(r \equiv 1 \pmod{8}\) with \(r \mid W_p\)
and \(4U_p^2 \equiv 1 \pmod{r}\). Then
\(\operatorname{ord}(\alpha) = 8p\) in the split algebra.
\end{theorem}


\begin{proof}
With \(4U_p^2 \equiv 1\):
\(\alpha^{2p} = V_p U_p s\). Squaring:
\(\alpha^{4p} = 2(V_p U_p)^2\). From \(V_p^2 \equiv 8U_p^2 - 4
\equiv -2\) and \(U_p^2 \equiv 1/4\):
\((V_p U_p)^2 = (-2)(1/4) = -1/2\), so
\(\alpha^{4p} = -1\). Hence \(\operatorname{ord}(\alpha) \mid 8p\)
but \(\operatorname{ord}(\alpha) \nmid 4p\); the only such divisors
are \(8\) and \(8p\). If \(\operatorname{ord}(\alpha) = 8\), then
\(\alpha^8 = (1+\sqrt 2)^8 = 577 + 408\sqrt 2 \equiv 1 \pmod r\)
would force \(r \mid \gcd(576, 408) = 24\), contradicting \(r \geq 17\).
Hence \(\operatorname{ord}(\alpha) = 8p\).\end{proof}


\subsection{Pell rank and the equivalence with Conjecture~\ref{conj:R3}}


\begin{definition}
The \emph{Pell rank of apparition} \(\rho(r)\) is the
smallest positive \(n\) with \(r \mid U_n\). We say \(r\) is a
\emph{primitive divisor} of \(U_n\) if \(\rho(r) = n\). Primitive divisors
exist for all \(n \geq 2\) \cite{ref2}.
\end{definition}


\begin{proposition}[Rank equivalence]\label{prop:rank-equivalence}
Let \(r \equiv 1 \pmod{8}\)
be a prime with \(r \mid W_p\). Suppose Condition (II) holds at \(r\)
in a composite \(W_p\), so that \(\operatorname{ord}(\alpha_0) = 4d\)
with \(d\) odd (where \(\alpha_0\) is the image of \(\omega_3\) in
one component of the split algebra, as in \S5.1). Write
\(\beta = 1+s \in \mathbb F_r\), so that \(\beta^2 = \alpha_0\); then
\(\operatorname{ord}(\beta) = 8d\). Then:
\begin{enumerate}
\item[\textup{(a)}] \(\rho(r) = 4d\), and \(r\) is a primitive divisor of \(U_{4d}\).
\item[\textup{(b)}] Conjecture~\ref{conj:R3} at \(r\) holds if and only if \(p \mid \rho(r)\).
\end{enumerate}
\end{proposition}

\begin{proof}
Justify \(\operatorname{ord}(\beta) = 8d\) first. Since
\(\beta^2 = \alpha_0\), \(\operatorname{ord}(\alpha_0) =
\operatorname{ord}(\beta)/\gcd(\operatorname{ord}(\beta),2)\). An
odd \(\operatorname{ord}(\beta)\) would force
\(\operatorname{ord}(\alpha_0)\) odd, contradicting
\(\operatorname{ord}(\alpha_0) = 4d\); so
\(\operatorname{ord}(\beta) = 2\operatorname{ord}(\alpha_0) = 8d\).

(a) From \(\operatorname{ord}(\beta) = 8d\) we have
\(\beta^{4d} = -1\). Squaring the identity
\(\alpha^n = V_n/2 + U_n\sqrt{2}\) and reducing modulo \(r\) with
\(s\) in place of \(\sqrt 2\) gives
\[
\beta^{2n} = (4U_n^2 + (-1)^n) + V_n U_n s,
\]
so \(\beta^{2n} = (-1)^n\) if and only if \(r \mid U_n\) (the
\(s\)-coefficient vanishes iff \(r \mid V_n\) or \(r \mid U_n\), and
the constant-term condition \(4U_n^2 \equiv 0 \pmod r\) forces
\(r \mid U_n\)).

For \(n = 4d\): \(\beta^{8d} = 1 = (-1)^{4d}\), so \(r \mid U_{4d}\).
For any proper divisor \(n\) of \(4d\): if \(n\) is even,
\(\beta^{2n} = 1\) requires \(8d \mid 2n\), so \(4d \mid n\),
impossible. If \(n\) is odd, \(\beta^{2n} = -1\) requires
\(\beta^{2n} = \beta^{4d}\), i.e., \(8d \mid (2n - 4d)\), i.e.,
\(4d \mid 2n - 4d\), i.e., \(2d \mid n\), impossible for odd \(n\).
So \(\rho(r) = 4d\) and \(r\) is a primitive divisor of \(U_{4d}\).

(b) Conjecture~\ref{conj:R3} says \(p \mid \operatorname{ord}(\beta) = 8d\).
Since \(p\) is odd, this holds iff \(p \mid d\), iff
\(p \mid 4d = \rho(r)\).
\end{proof}
\begin{definition}[Compatible triple]\label{def:compatible-triple}
A \emph{compatible triple} \((d, r, p)\) consists of:
\begin{itemize}
\item an odd integer \(d > 1\);
\item a prime \(r \equiv 1 \pmod{8}\) that is a primitive divisor of
  \(U_{4d}\);
\item a prime \(p \geq 5\) with \(\operatorname{ord}_r(2) = 2p\) and
  \(d \mid W_{p-2}\).
\end{itemize}
\end{definition}


By Proposition~\ref{prop:rank-equivalence}, if Condition (II) passes at some split factor of
a composite \(W_p\), then a compatible triple exists. The divisibility
constraint \(d \mid W_{p-2}\) relates to the general theory of
Sgobba \cite{ref16}.

\subsection{Quartic residuosity of 2}


\begin{theorem}[Quartic residuosity at primitive Pell divisors]\label{thm:quartic-residuosity}
Let \(n \geq 1\) be odd. If \(r\) is a prime with
\(r \equiv 1 \pmod{8}\) that is a primitive divisor of \(U_{4n}\),
then \(2\) is a fourth-power residue modulo \(r\):
\[
2^{(r-1)/4} \equiv 1 \pmod{r}.
\]
\end{theorem}


\begin{proof}
In \(\mathbb{Q}(\zeta_8) = \mathbb{Q}(i, \sqrt{2})\):
\[
\sqrt{2} = (1 - \zeta_8)^2 \cdot \zeta_8^3 \cdot (1 + \sqrt{2}).
\tag{$\ast$}
\]
Indeed, \((1-\zeta_8)^2 = (1+i)(1-\sqrt{2})\),
so \((1-\zeta_8)^2 \cdot \zeta_8^3 = 2 - \sqrt{2}\), and
\((2 - \sqrt{2})(1+\sqrt{2}) = \sqrt{2}\), confirming~\((\ast)\).

Since \(r \equiv 1 \pmod{8}\), reduce \((\ast)\) modulo a prime
\(\mathfrak{P} \mid r\) in \(\mathbb{Z}[\zeta_8]\), sending
\(\zeta_8 \mapsto \omega\), \(\sqrt{2} \mapsto s\),
\(1+\sqrt{2} \mapsto \alpha\):
\[
s = (1-\omega)^2 \cdot \omega^3 \cdot \alpha \quad\text{in }
\mathbb{F}_r.  \tag{$\ast\ast$}
\]

Raise to \((r-1)/2\). Since \(r\) is a primitive divisor of
\(U_{4n}\): \(\operatorname{ord}(\alpha) = 8n\) and
\(\alpha^{4n} = -1\). Write \(r - 1 = 8nM\).

\begin{itemize}
\item \([(1-\omega)^2]^{(r-1)/2} = 1\) by Fermat.
\item \([\omega^3]^{(r-1)/2} = (-1)^{nM}\)
  (since \(\operatorname{ord}(\omega) = 8\)).
\item \(\alpha^{(r-1)/2} = (-1)^M\).
\end{itemize}

Therefore \(s^{(r-1)/2} = (-1)^{(n+1)M}\). Since \(n\) is odd,
\(n+1\) is even, giving \(s^{(r-1)/2} = 1\) and
\(2^{(r-1)/4} = s^{(r-1)/2} = 1\).
\end{proof}
\begin{remark}
The parity of \(n\) is essential: for even \(n\),
2 need not be a quartic residue at primitive divisors of \(U_{4n}\).
\end{remark}

\subsection{No Compatible Triples}\label{sec:nct}


\begin{theorem}[{NCT for $d \leq 81$}]\label{thm:NCT-81}
For each odd \(d\) with
\(1 \leq d \leq 81\), no compatible triple \((d, r, p)\) exists.
\end{theorem}


\begin{proof}
Call \(d\) \emph{admissible} if every prime factor \(q\) of \(d\)
satisfies \(v_2(\operatorname{ord}_q(2)) = 1\); otherwise
\(q \mid W_{p-2}\) is impossible (if \(\operatorname{ord}_q(2)\) is
odd, no half-order \(m_q\) exists; if
\(4 \mid \operatorname{ord}_q(2)\), then \(m_q\) is even but
\(p - 2\) is odd). This parity obstruction eliminates 80\% of odd
\(d \leq 500\); within \([1, 81]\) it leaves 12 values.

For \(d = 1, 3\): \(U_4 = 12\) and \(U_{12} = 13860\) have no prime
factor \(\equiv 1 \pmod{8}\), so no split \(r\) exists. For the
remaining 10 values \(d \in \{9, 11, 19, 27, 33, 43, 57, 59,
67, 81\}\): \(U_{4d}\) is completely factored (the largest,
\(U_{324}\), has 124 digits). Every primitive prime divisor
\(r \equiv 1 \pmod{8}\) of \(U_{4d}\) is checked (18 such primes
across the 10 values of \(d\); the largest is
\(r = 306{,}416{,}447{,}870{,}523{,}097\) for \(d = 59\)): in all cases,
\(\operatorname{ord}_r(2)\) is either odd or
\(\operatorname{ord}_r(2)/2\) is composite, so no compatible triple
exists. The data for each admissible \(d\):

\begin{center}
\small
\begin{tabular}{rrl}
\hline
\(d\) & \(\#\{r \equiv 1\,(8) \text{ prim.}\}\) & Obstruction \\
\hline
9 & 1 & \(\operatorname{ord}_{73}(2) = 9\) (odd) \\
11 & 1 & \(\operatorname{ord}_{89}(2) = 11\) (odd) \\
19 & 2 & odd or composite \(\operatorname{ord}/2\) \\
27 & 2 & composite \(\operatorname{ord}/2\) \\
33 & 1 & \(\operatorname{ord}_{16633}(2) = 2079\) (odd) \\
43 & 2 & composite \(\operatorname{ord}/2\) \\
57 & 2 & odd or composite \(\operatorname{ord}/2\) \\
59 & 2 & odd or composite \(\operatorname{ord}/2\) \\
67 & 3 & composite \(\operatorname{ord}/2\) \\
81 & 2 & composite \(\operatorname{ord}/2\) \\
\hline
\end{tabular}
\end{center}

\noindent Complete factorizations of each \(U_{4d}\) are in the Zenodo
archive.
\end{proof}
Combined with Theorems~\ref{thm:case-A}--\ref{thm:case-4U},
the three explicit cases and Theorem~\ref{thm:NCT-81} provide an
unconditional proof whenever the rank parameter \(d\) is admissible
(odd) and \(d \leq 81\). Among the 185 known split factors,
Cases A, B, and \(4U_p^2 \equiv 1\) cover 29\%
(Theorems~\ref{thm:case-A}--\ref{thm:case-4U}); most remaining
Case C factors have \(d > 200\) and are covered by
computational R3 verification (\S9.2).

\begin{proposition}[{$q$-th power necessity}]\label{prop:qth-power}
In any compatible
triple \((d, r, p)\), for every prime \(q \mid d\) with
\(q \equiv 3 \pmod{8}\): 2 must be a \(q\)-th power residue
modulo \(r\).
\end{proposition}


\begin{proof}
Suppose 2 is not a \(q\)-th power mod \(r\). Then
\(q \mid \operatorname{ord}_r(2) = 2p\), so \(q = p\). Then
\(p \mid d \mid W_{p-2}\), contradicting Lemma~\ref{lem:p-nmid-Wp2}.\end{proof}


\begin{conjecture}[No Compatible Triples, general]\label{conj:nct-full}
For every odd \(d > 1\), no compatible triple
\((d, r, p)\) exists.
\end{conjecture}


\begin{remark}\label{rem:qpro-false}
A natural strengthening of Proposition~\ref{prop:qth-power} --- that
\(v_2(\operatorname{ord}_r(2)) \geq 2\) whenever 2 is a \(q\)-th
power modulo a primitive divisor \(r\) of \(U_{4q}\) --- is false.
For \(q = 163\), the primitive divisor
\(r = 12036824328615957881\) of \(U_{652}\) satisfies
\(2^{(r-1)/163} \equiv 1 \pmod{r}\), yet
\(v_2(\operatorname{ord}_r(2)) = 1\).
Conjecture~\ref{conj:nct-full} is unaffected: \(\operatorname{ord}_r(2)/2\) is
composite (\S9.3).
\end{remark}

\begin{remark}[Structural source of the split-branch gap]\label{rem:split-gap}
All \(\mathbb{F}_r\)-internal algebraic identities are tautological
(confirmed across the 185 known split factors). The quartic
residuosity of 2 is proved unconditionally (Theorem~\ref{thm:quartic-residuosity}), but
higher-power conditions do not yield further obstructions
(Remark~\ref{rem:qpro-false}). NCT should be attacked directly via compatible
triples.
\end{remark}


\section{The \texorpdfstring{\(r \equiv 3 \pmod{8}\)}{r == 3 (mod 8)} branch}


\subsection{The gap structure}


For \(r \equiv 3 \pmod{8}\) with \(r \mid W_p\) (composite):
\((2/r) = -1\), so \(\mathbb{Z}[\sqrt{2}]/(r) \cong \mathbb{F}_{r^2}\).
The element \(\omega_3\) has norm 1 and therefore lies in the norm-1
subgroup \(\mu_{r+1}\); by Lemma~\ref{lem:v2-classification},
\(v_2(\operatorname{ord}(\omega_3)) = 2\), so
\(\operatorname{ord}(\omega_3) = 4g\) for some odd \(g\). Combined
with \(\operatorname{ord}(\omega_3) \mid r+1\) and
\(v_2(r+1) = 2\), this gives \(g \mid (r+1)/4\).

\textbf{Descent of Condition (II).} Condition (II) at \(r\) is
\(\omega_3^{(N+1)/2} = -1\), i.e., \(\omega_3^{2W_{p-2}} = -1\)
(since \(N+1 = 4W_{p-2}\)). With \(\operatorname{ord}(\omega_3) = 4g\):
\(4g \mid 4W_{p-2}\) iff \(g \mid W_{p-2}\), and \(4g \nmid 2W_{p-2}\)
since \(W_{p-2}\) is odd. Therefore Condition (II) passes at \(r\)
if and only if \(g \mid W_{p-2}\).

\textbf{Notation.} For inert \(r \mid W_p\), define
\[
G_r = \gcd\!\left(\frac{r+1}{4},\, W_{p-2}\right).
\]
If Condition (II) passes at \(r\), then \(g \mid G_r\).

\begin{remark}\label{rem:frobenius-norm}
Condition (II) already implies
\(\omega_3^N \equiv \omega_3^{-1} \pmod{N}\), so quadratic Frobenius
strengthenings \cite{ref8} add nothing. The norm
\(\operatorname{Norm}(1+\sqrt{2}) = -1\) constrains only the 2-part
of \(\operatorname{ord}(\omega_3)\); the odd part \(g\) is not
constrained by the norm argument alone.
\end{remark}


\subsection{Parity and valuation obstructions}


We classify odd primes \(q\) by \(v_2(\operatorname{ord}_q(2))\):
\emph{Class~I} if \(\operatorname{ord}_q(2)\) is odd (\(v_2 = 0\));
\emph{Class~II} if \(v_2(\operatorname{ord}_q(2)) \geq 2\);
\emph{Class~III} if \(v_2(\operatorname{ord}_q(2)) = 1\). Only
Class~III primes can divide \(W_n\) for odd \(n\), since
\(q \mid W_n\) requires \(2^n \equiv -1 \pmod{q}\), hence
\(\operatorname{ord}_q(2) = 2m\) with \(m\) odd. Define
\(A = \{q \text{ prime} : v_2(\operatorname{ord}_q(2)) = 1\}\),
the set of Class~III primes.
Every prime \(q \equiv 3 \pmod{8}\) lies in \(A\); primes
\(q \equiv 5\) or \(7 \pmod{8}\) do not (they are Class~II and
Class~I respectively, by the argument in Lemma~\ref{lem:mod-8-exclusion}).

\begin{theorem}[Parity obstruction for inert factors]\label{thm:parity-obstruction}
Let
\(r \equiv 3 \pmod{8}\) with \(r \mid W_p\) (composite) and
\(g = \operatorname{ord}_r(\omega_3)/4\). If \(g\) has a prime factor
\(q \notin A\), then \(g \nmid W_{p-2}\) and Condition (II) fails.
\end{theorem}


\begin{proof}
If \(q \notin A\), then \(v_2(\operatorname{ord}_q(2)) \neq 1\),
so \(q \nmid W_n\) for any odd \(n\) (the condition \(q \mid W_n\)
requires \(\operatorname{ord}_q(2) = 2m_q\) with \(m_q\) odd, hence
\(v_2(\operatorname{ord}_q(2)) = 1\)). Since \(q \mid g\) and
\(q \nmid W_{p-2}\): \(g \nmid W_{p-2}\).\end{proof}


\begin{theorem}[Generalized LTE for Class III primes]\label{thm:generalized-LTE}
Let
\(q\) be a prime with \(v_2(\operatorname{ord}_q(2)) = 1\), say
\(\operatorname{ord}_q(2) = 2m\) with \(m\) odd. For odd \(n\)
with \(m \mid n\):
\[
v_q(2^n + 1) = v_q(2^m + 1) + v_q(n/m).
\]
For \(q \geq 5\): \(v_q(W_n) = v_q(2^m + 1) + v_q(n/m)\).
For \(q = 3\): \(v_3(W_n) = v_3(n)\).
\end{theorem}


\begin{proof}
Write \(n = m \cdot \ell\) with \(\ell\) odd. Then
\(2^n + 1 = (2^m)^\ell + 1\). The lifting-the-exponent lemma
with \(a = 2^m\), \(b = 1\) (so that \(q \mid a + b\)),
exponent \(\ell\) (odd), gives
\(v_q(2^n + 1) = v_q(2^m + 1) + v_q(\ell)\).
Since \(3W_n = 2^n + 1\): for \(q \geq 5\),
\(v_q(W_n) = v_q(2^n + 1)\). For \(q = 3\): \(m = 1\),
\(v_3(2^n+1) = 1 + v_3(n)\), and
\(v_3(W_n) = v_3(2^n+1) - 1 = v_3(n)\).\end{proof}


Computationally, \(v_q(2^m + 1) = 1\) for all 194 Class III primes
\(q < 5000\), where \(m = \operatorname{ord}_q(2)/2\); i.e., no
\emph{Wieferich-at-minus-1} behaviour (\(q^2 \mid 2^m + 1\)) is observed.
This is the multiplicative-order-\(2m\) analogue of the Wieferich
condition \(q^2 \mid 2^{q-1} - 1\).

\textbf{Valuation obstruction.} If \(v_q(g) > v_q(W_{p-2})\) for some
\(q \in A\) with \(q \mid g\), then \(g \nmid W_{p-2}\) and
Condition (II) fails. For \(q = 3\): this occurs whenever
\(v_3(g) > v_3(p-2)\).

Among 869 inert factors of composite \(W_p\) with \(r < 2 \times 10^5\):
approximately 86\% are blocked by parity or valuation obstructions
(\S9.5).

\subsection{Pell-sequence exclusion}


The following result gives a direct divisibility constraint on
inert factors.

\begin{theorem}[Pell-sequence divisibility]\label{thm:pell-divisibility}
Let
\(r \equiv 3 \pmod{8}\) be a prime with \(r \mid W_p\), and suppose
\(\operatorname{ord}_r(\omega_3) = 4d\) in \(\mathbb{F}_{r^2}\). Then
\[
r \mid \gcd\!\left(\frac{V_{4d}+2}{2},\; U_{4d}\right).
\]
\end{theorem}


\begin{proof}
The condition \(\operatorname{ord}_r(\omega_3) = 4d\) gives
\(\omega_3^{2d} = -1\) in \(\mathbb{F}_{r^2}\). Since
\(\omega_3 = \alpha^2\): \(\alpha^{4d} = -1\). Expanding
\(\alpha^{4d} = V_{4d}/2 + U_{4d}\sqrt{2}\) and comparing with
\(-1 = -1 + 0 \cdot \sqrt{2}\):
\(r \mid U_{4d}\) and \(r \mid (V_{4d}/2 + 1) = (V_{4d}+2)/2\).
(Since \(r\) is inert, \(\mathbb{F}_{r^2} = \mathbb{F}_r(\sqrt{2})\)
and \(\{1, \sqrt{2}\}\) is an \(\mathbb{F}_r\)-basis, so coefficient
comparison is valid. Note: \(V_{4d}\) is even, so \((V_{4d}+2)/2\)
is an integer.)\end{proof}


\begin{definition}[Pell exclusion number]\label{def:pell-exclusion-number}
For odd \(d \geq 1\), define the \emph{Pell exclusion number}
\[
N_d = \gcd\!\left(\frac{V_{4d}+2}{2},\; U_{4d}\right).
\]
\end{definition}


\begin{proposition}[Closed form for $N_d$]\label{prop:Nd-closed-form}
For odd \(d \geq 1\),
\[
N_d \;=\; V_{2d} \;=\; V_d^{\,2} + 2.
\]
\end{proposition}

\begin{proof}
Three standard Pell identities:
(i)~\(V_{4d} = V_{2d}^{\,2} - 2(-1)^{2d} = V_{2d}^{\,2} - 2\), so
\(V_{4d} + 2 = V_{2d}^{\,2}\);
(ii)~\(U_{4d} = U_{2d}\, V_{2d}\);
(iii)~\((V_{2d}/2)^2 - 2\,U_{2d}^{\,2} = (-1)^{2d} = 1\), so
\(\gcd(V_{2d}/2,\, U_{2d}) = 1\).
Since \(V_{2d}\) is even (write \(V_{2d} = 2k\)),
\(N_d = \gcd(2k^2,\; 2k\,U_{2d}) = 2k\cdot\gcd(k, U_{2d}) = 2k = V_{2d}\).
Finally, \(V_{2d} = V_d^{\,2} - 2(-1)^d = V_d^{\,2} + 2\) since \(d\) is odd.
\end{proof}

\begin{remark}\label{rem:Nd-norm}
Since \(N_d = V_d^{\,2} + 2 = \operatorname{N}_{\mathbb{Z}[\sqrt{-2}]}(V_d + \sqrt{-2})\),
every prime \(r \equiv 3 \pmod{8}\) dividing \(N_d\) splits in
\(\mathbb{Z}[\sqrt{-2}]\) (because \((-2/r) = 1\) for such \(r\)).
\end{remark}


\begin{theorem}[{Pell-sequence exclusion for $d \leq 43$}]\label{thm:pell-exclusion}
For each odd \(d\) in the admissible set
\(\{1, 3, 9, 11, 19, 27, 33, 43\}\) (those \(d \leq 43\) with
all prime factors in \(A\)): no prime \(r \equiv 3 \pmod{8}\)
dividing a composite \(W_p\) can have
\(\operatorname{ord}_r(\omega_3) = 4d\).
\end{theorem}


\begin{proof}
By Theorem~\ref{thm:pell-divisibility}, such \(r\) would divide
\(N_d = V_{2d}\) (Proposition~\ref{prop:Nd-closed-form}).

For \(d = 1\): \(N_1 = V_1^2 + 2 = 6 = 2 \cdot 3\). The only prime
\(\equiv 3 \pmod{8}\) is 3. Since \(\gcd(W_p, 3) = 1\) for
\(p \geq 5\): no valid \(r\) exists.

For \(d = 3\): \(N_3 = V_3^2 + 2 = 14^2 + 2 = 198 = 2 \cdot 3^2 \cdot 11\). Primes
\(\equiv 3 \pmod{8}\): 3 and 11. We have \(3 \nmid W_p\). For
\(r = 11\): \(\operatorname{ord}_{11}(2) = 10\), giving \(p = 5\),
but \(W_5 = 11\) is prime. So 11 does not divide any composite \(W_p\).

For \(d \in \{9, 11, 19, 27, 33, 43\}\): \(N_d = V_d^2 + 2\) is
completely factored (the largest, \(N_{43}\), has 33 digits). Each
prime factor \(r \equiv 3 \pmod{8}\) with \(r > 11\) is checked:
either \(\operatorname{ord}_r(2)\) has \(v_2 \neq 1\) (so
\(r \nmid W_p\) for any prime \(p\)), or
\(\operatorname{ord}_r(2)/2\) is composite (so no prime \(p\)
satisfies \(\operatorname{ord}_r(2) = 2p\)), or the unique prime
\(p = \operatorname{ord}_r(2)/2\) gives \(d \nmid W_{p-2}\).
In all cases, no valid configuration exists.

For \(d\) not in the admissible set (i.e., \(d \leq 43\) with some
prime factor \(q \notin A\)): Theorem~\ref{thm:parity-obstruction} gives
\(g \nmid W_{p-2}\), so \(\operatorname{ord}_r(\omega_3)/4 \neq d\)
in any Condition-(II)-passing scenario.
\end{proof}
\begin{corollary}[{Inert factors with small $G_r$}]\label{cor:small-Gr}
Let
\(r \equiv 3 \pmod{8}\) divide a composite \(W_p\). If
\(G_r \leq 43\), then Condition (II) fails at \(r\).
\end{corollary}


\begin{proof}
If Condition (II) passes, then
\(g = \operatorname{ord}_r(\omega_3)/4\) divides \(G_r\), so
\(g \leq 43\). Since \(g\) is odd and \(g \mid W_{p-2}\), all prime
factors of \(g\) lie in \(A\) (else Theorem~\ref{thm:parity-obstruction} blocks). So \(g\) is
in the admissible set of Theorem~\ref{thm:pell-exclusion}. But Theorem~\ref{thm:pell-exclusion} says no valid
\(r\) exists for such \(g\). Contradiction.\end{proof}


\subsection{Bounds on the order parameter}


\begin{lemma}\label{lem:class-III-div}
For \(q \in A\) with \(q > 3\) and \(m_q =
\operatorname{ord}_q(2)/2\): \(m_q \mid (p-2)\) if and only if
\(q \mid W_{p-2}\).
\end{lemma}


\begin{proof}
Since \(v_2(\operatorname{ord}_q(2)) = 1\):
\(2^{m_q} \equiv -1 \pmod{q}\). For odd \(n\) with \(m_q \mid n\):
\(2^n = (-1)^{n/m_q} = -1\), so \(q \mid 3W_n\) and \(q \mid W_n\).
Conversely, \(q \mid W_n\) gives \(2^n \equiv -1\), so
\(m_q \mid n\).\end{proof}


\begin{lemma}[Lower bound on $g$]\label{lem:lower-bound-g}
For inert \(r \mid W_p\) (composite):
\(g > \log(2p)/\log(3+2\sqrt{2}) \approx 0.393 \cdot \log_2(p)\).
\end{lemma}

\begin{proof}
From \(\operatorname{ord}_r(\omega_3) = 4g\) we get
\(\omega_3^{2g} \equiv -1 \pmod r\), hence
\(V_{2g} \equiv \omega_3^g + \omega_3^{-g} \equiv \omega_3^g +
(-\omega_3^g) \equiv 0 \pmod r\), i.e.\ \(r \mid V_{2g}\).
Since \(r \mid W_p\) is inert, \(\operatorname{ord}_r(2) = 2p\)
divides \(r-1\), so \(r \geq 2p+1\). Combining with
\(V_{2g} = \omega_3^g + \omega_3^{-g} < (3+2\sqrt{2})^g + 1\):
\(2p+1 \leq r \leq V_{2g} < (3+2\sqrt{2})^g + 1\), so
\(g > \log(2p)/\log(3+2\sqrt{2})\).
\end{proof}


\textbf{Product bound.} For a given \(p\):
\[
G_{\max}(p) = 3^{v_3(p-2)} \cdot \prod_{\substack{q > 3,\; q \equiv 3\,(8) \\
m_q \mid (p-2)}} q^{v_q(W_{p-2})}
\]
is an upper bound on any compatible \(g\). When
\(G_{\max}(p) < 3\) (e.g., when \(p - 2\) is prime), no inert
factor can satisfy Condition (II).

\subsection{Coverage and computational status}


Corollary~\ref{cor:small-Gr} closes the inert branch unconditionally whenever
\(G_r \leq 43\), covering 845 of the 869 inert factors structurally
analysed in \S9.5. A natural subclass of the residual configurations
(\(G_r > 43\)) --- those inert primes for which \((r+1)/4\) itself
has every odd prime factor in \(A\) --- has natural density zero
within the primes \(r \equiv 3 \pmod 8\) (Proposition~\ref{prop:density-zero},
unconditional). The full residual set is slightly larger, and the
structural barrier to a \emph{pointwise} closure --- together with the
reason a naive GRH/Chebotarev density argument does not supply
one --- is discussed in \S10, Problem 2.

No composite \(W_p\) globally passes Condition~(II) in the
684{,}965{,}381-factor survey of \S9.5.

\subsection{Density and structural limits}

The norm \(N(\alpha) = -1\) constrains only the 2-part of orders,
and the odd part \(g\) is not constrained by the norm argument alone
(Remark~\ref{rem:frobenius-norm}). The generalized LTE (Theorem~\ref{thm:generalized-LTE}) controls \(q\)-adic
valuations of \(W_{p-2}\) for primes \(q \in A\), and the
Pell-sequence exclusion constrains the \emph{factor} \(r\) rather than
the \emph{configuration}, succeeding for all \(g \leq 43\) because
\(N_d = V_d^2 + 2\) (Proposition~\ref{prop:Nd-closed-form}) is small enough to factor completely at those \(d\). Beyond
\(d = 43\), sporadic large prime factors of \(N_d\) survive all
local checks. However, the surviving configurations are negligible
in density:

\begin{proposition}[Density zero for pure Class III inert factors]\label{prop:density-zero}
The set of primes \(r \equiv 3 \pmod{8}\) for which every odd prime
factor of \((r+1)/4\) belongs to Class III has density zero within
the primes \(r \equiv 3 \pmod{8}\).
\end{proposition}

\begin{proof}
Let \(S = \{q \text{ prime} : v_2(\operatorname{ord}_q(2))
\geq 2\}\). Every prime \(q \equiv 5 \pmod{8}\) satisfies
\(v_2(\operatorname{ord}_q(2)) = 2\) (by the argument in
Lemma~\ref{lem:mod-8-exclusion}), so \(S\) contains the full density-\(1/4\) set
of primes \(\equiv 5 \pmod{8}\); in particular,
\(\sum_{q \in S} 1/q = \infty\).

If every odd prime factor of \((r+1)/4\) is Class III, then
\(q \nmid (r+1)\) for all \(q \in S\) (since each \(q \in S\)
is odd and not Class III). Equivalently,
\(r \not\equiv -1 \pmod{q}\) for every \(q \in S\).

Fix a finite subset \(S' \subset S\) and let
\(M = 8 \prod_{q \in S'} q\). The conditions
\(r \equiv 3 \pmod{8}\) and \(r \not\equiv -1 \pmod{q}\)
for each \(q \in S'\) determine a union of residue classes
modulo \(M\) (by CRT, since the moduli \(8, q_1, q_2, \ldots\) are
pairwise coprime). By the prime number theorem for arithmetic
progressions with the fixed modulus \(M\), the natural density
of primes in this union, relative to all primes
\(r \equiv 3 \pmod{8}\), is
\[
\prod_{q \in S'} \frac{q - 2}{q - 1}.
\]
(Among the \(q-1\) nonzero residue classes modulo \(q\), exactly
one --- namely \(-1\) --- is forbidden; the remaining
\(q-2\) classes are allowed, giving the factor \((q-2)/(q-1)\).)

Since every prime \(r\) in our target set satisfies
\(q \nmid (r+1)\) for all \(q \in S\), it satisfies these
conditions for every finite \(S' \subset S\). Therefore the
upper relative density of the target set among primes
\(r \equiv 3 \pmod{8}\) is at most
\(\prod_{q \in S'} (q-2)/(q-1)\) for every finite
\(S' \subset S\). From \(\sum_{q \in S} 1/q = \infty\) and
\(1/(q-1) \geq 1/q\) for all \(q \geq 2\), we have
\(\sum_{q \in S} 1/(q-1) = \infty\); hence the infimum of the
products \(\prod_{q \in S'} (1 - 1/(q-1))\) over finite
\(S' \subset S\) is zero, and the (relative) density is zero.\end{proof}

Note that Proposition~\ref{prop:density-zero} applies to the natural subclass in which
\((r+1)/4\) itself has only Class III prime factors; this is a
proper subset of the full residual configurations (which require
only that the order parameter \(g = \operatorname{ord}_r(\omega_3)/4\)
has every prime factor in \(A\), not that all of \((r+1)/4\) does).
Even for this subclass, the density-zero bound is not an \emph{eventual}
(pointwise) closure: the exceptional set, though thin, is infinite.

\begin{remark}[Naive model for the inert branch]\label{rem:inert-heuristic}
For comparison, the split branch (\S9.8, footnote) admits a naive
independence model with expected failure count
\(E = \sum 1/p \approx 0.76\), so that zero failures among
\({\sim}14.8 \times 10^6\) split factors is consistent with, but
not strong evidence for, Conjecture~\ref{conj:R3}. For the inert
branch, no analogous convergent model exists: the sum of
local-pass probabilities over the tested inert factors diverges,
because the algebraic obstructions (parity, Pell exclusion) that
block most configurations are not captured by a uniform random
model. The 684{,}965{,}381 zero-failure count reflects the
algebraic structure of \S\S5--6 rather than a small Poisson
expectation; the structure is proved to suffice for \(d \leq 43\)
but remains open beyond.
\end{remark}

A naive GRH/Chebotarev density argument in the style of Hooley
\cite{ref10} and Heath-Brown \cite{ref9} does \emph{not} supply the missing closure
for the full residual set either: the relevant Hooley-type density
\(\delta = \prod_{q \equiv 3\,(8),\,q \in A,\,m_q \mid (p-2)}
(1 - 1/q)\)
is positive (approximately \(0.66\) for all \(p\) in the tested
range), reflecting the fact that the constraints ``\(q \in A\)''
and ``\(m_q \mid (p-2)\)'' are mutually compatible with a positive
density of primes \(r\). An unconditional closure would require
one of two ingredients: either that every inert \(g\) contains a
prime \(q \notin A\), or a valuation excess via Theorem~\ref{thm:generalized-LTE}. Both
routes face Chebotarev-type obstructions analogous to those of
Artin's primitive root conjecture (see \cite{ref15} for a survey). Unlike
Artin's conjecture proper, however, the problems here carry
additional algebraic structure --- Pell-sequence divisibility for
NCT, and LTE valuation control for inert factors --- that may
admit approaches beyond the Chebotarev density theorem.

\subsection{Multi-factor reduction}\label{sec:multi-factor}

The preceding single-factor analysis leaves the residual gap
\(d > 43\) open factor by factor. We now combine global constraints
to show that any all-inert composite counterexample requires a
multi-factor configuration that is tightly constrained.

\begin{definition}[Danger triple]\label{def:danger-triple}
A \emph{danger triple} \((d, r, p)\) consists of an odd integer
\(d > 43\) with every prime factor in~\(A\), a prime
\(r \equiv 3 \pmod{8}\) that is a primitive divisor of \(V_{2d}\)
with \(\operatorname{ord}_r(\omega_3) = 4d\) in
\(\mathbb{F}_{r^2}\), and the unique prime \(p\) with
\(\operatorname{ord}_r(2) = 2p\), satisfying \(d \mid W_{p-2}\).
A danger triple is \emph{phantom} if \(W_p\) is prime.
\end{definition}

By Theorem~\ref{thm:pell-divisibility}, any inert factor \(r\) of
a composite \(W_p\) satisfying Condition~(II) with order parameter
\(d > 43\) gives rise to a danger triple \((d, r, p)\).

\begin{theorem}[Multi-factor reduction]\label{thm:multi-factor}
Let \(W_p\) be composite with every prime factor
\(r \equiv 3 \pmod{8}\). If Condition~(II) holds at every factor,
then \(W_p\) has \(t \geq 3\) prime factors, and:
\begin{enumerate}
\item[\textup{(a)}] each factor \(r_i\) has order parameter
  \(d_i > 43\) with every prime factor of \(d_i\) in \(A\);
\item[\textup{(b)}] at most two factors have
  \(d_i \in \{57, 67, 331\}\);
\item[\textup{(c)}] at least one factor has
  \(d_i \notin \{57, 67, 331, 683\}\).
\end{enumerate}
\end{theorem}

\begin{proof}
\emph{Three or more factors.}\; Since each
\(r_j \equiv 3 \pmod{8}\):
\(\prod r_j \equiv 3^t \pmod{8}\). Because
\(W_p \equiv 3 \pmod{8}\) (Lemma~\ref{lem:mod-8-exclusion}), \(t\) is odd;
as \(W_p\) is composite, \(t \geq 3\).

\emph{Part~(a)} follows from Theorem~\ref{thm:pell-exclusion}
(\(d_i > 43\)) and Theorem~\ref{thm:parity-obstruction} (prime
factors of \(d_i\) in \(A\)).

\emph{Part~(b).}\; Complete factorisation of \(V_{114}\) and
\(V_{134}\) enumerates all primitive inert divisors for \(d = 57\)
and \(d = 67\). The danger \(p\)\nobreakdash-sets are
\(\{11, 85684865933\}\) and \(\{24822855876938710967\}\)
respectively. These sets are disjoint, and all factors of a
given \(W_p\) share the same~\(p\); thus no single \(p\) can
simultaneously produce a factor with \(d = 57\) and another with
\(d = 67\), giving at most one factor from
\(\{57, 67\}\). For \(d = 331\): the primitive part of \(V_{662}\)
(253 digits) has a 238-digit unfactored cofactor; among the
identified primitive divisors, one danger triple arises
(\(p = 17\), phantom). Regardless of how many further danger
triples may lie in the cofactor, at most one factor of any
given \(W_p\) can have \(d_i = 331\), since a single \(p\) selects
at most one primitive divisor per \(d\)-value. Combined: at most
two factors with \(d_i \in \{57, 67, 331\}\).

\emph{Part~(c).}\; Since \(t \geq 3\) and at most two factors
have \(d_i \in \{57, 67, 331\}\), at least one factor has \(d_i\)
outside this set. A primitive-divisor survey of all 79
admissible \(d \leq 1000\) (those whose prime factors all
satisfy \(v_2(\operatorname{ord}_q(2)) = 1\); the parity
obstruction eliminates the rest; enumeration in the
supplementary data) identifies exactly one further danger
triple, at \(d = 683\) (\(p = 13\), \(r = 2731 = W_{13}\),
phantom). Since \(W_{13}\) is prime, \(d = 683\) cannot produce a
factor of a composite \(W_p\). No danger triple is identified at
any other admissible \(d \leq 1000\). Hence \(d_i > 1000\) for
at least one factor.\footnote{The survey factors the primitive
part of \(V_{2d}\) as far as practical (complete for
\(d \in \{57, 67\}\), partial for larger \(d\) including \(d = 331\));
no unidentified danger triple in an unfactored cofactor can
affect the bound of Part~(b), which relies on same-\(d\) exclusion
rather than completeness of the factorisation. The conclusion
\(d_i > 1000\) is the empirical finding of the survey.}
\end{proof}

\begin{remark}\label{rem:phantom-structure}
Of the five identified danger triples, three are phantom:
\((d, r, p) = (57, 683, 11)\) with \(W_{11} = 683\) prime,
\((331, 43691, 17)\) with \(W_{17} = 43691\) prime, and
\((683, 2731, 13)\) with \(W_{13} = 2731\) prime. The two real
dangers --- at \((d, p) = (57, 85684865933)\) and
\((67, 24822855876938710967)\) --- are each disposed of at
secondary factors (\S9.8).
\end{remark}


\section{Summary of the case analysis}


The factor-by-factor status is recorded in the table of
\S\ref{sec:status-table}. Two structurally independent problems
remain open: NCT for general~\(d\) (\S\ref{sec:nct}), a finite
local verification at each~\(d\) (proved through \(d = 81\));
and the universal order-parameter bound for inert factors
(\S6.4), currently closed only in density
(Proposition~\ref{prop:density-zero}).

Section~\ref{sec:pinning-reductions} further reduces the inert
branch by a sequence of cross-case unconditional arguments
(Order-Pinning Lemma, Multi-Factor Pinning, Phantom Exclusion of
the finite \(d \leq 1000\) danger catalog, the Platinum Lemma at
\(r \leq 10^{12}\), and a deterministic Second-Moment Reduction)
to the single claim that no composite \(W_p\) carries two
distinct residual inert danger factors \(r_1, r_2 > 10^{12}\) at
the same~\(p\).



\section{Unconditional reductions from pinning}\label{sec:pinning-reductions}

This section records the unconditional reductions that act
uniformly across all candidate composite \(W_p\), beyond the
branch-by-branch analysis of \S\S5--6. None of the results here
depends on Conjecture~\ref{conj:R3} or on further computational
input at a single factor; each is a structural consequence of
\(r \mid W_p\) together with the Condition~(II) hypothesis. The
end product (Theorem~\ref{thm:second-moment-reduction}) reduces
the all-inert branch of Conjecture~\ref{conj:sufficiency} to the
single deterministic claim that no composite~\(W_p\) admits a
pair of distinct residual inert danger factors above~\(10^{12}\).

\subsection{Order-Pinning Lemma}

\begin{theorem}[Order-Pinning]\label{thm:order-pinning}
Let \(r\) be a prime factor of \(W_p\) with \(r > 3\). Then
\(\operatorname{ord}_r(2) = 2p\).
\end{theorem}

\begin{proof}
\(r \mid W_p\) means \(r \mid 2^p + 1\), so \(2^p \equiv -1
\pmod r\), giving \(2^{2p} \equiv 1 \pmod r\); hence
\(\operatorname{ord}_r(2) \mid 2p\) and
\(\operatorname{ord}_r(2) \nmid p\). Since \(p\) is prime, the
divisors of \(2p\) are \(\{1, 2, p, 2p\}\); the only divisors not
dividing \(p\) are \(\{2, 2p\}\). If
\(\operatorname{ord}_r(2) = 2\) then \(r \mid 3\), contradicting
\(r > 3\). Hence \(\operatorname{ord}_r(2) = 2p\).
\end{proof}

\subsection{Multi-Factor Pinning}

\begin{theorem}[Multi-Factor Pinning]\label{thm:multi-factor-pinning}
Let \(W_p\) be composite. All prime factors \(r > 3\) of \(W_p\)
share a common value \(\operatorname{ord}_r(2)/2 = p\).
Consequently, if two composite Wagstaff numbers \(W_{p_1}\) and
\(W_{p_2}\) with \(p_1 \neq p_2\) are given, they share no prime
factor \(r > 3\).
\end{theorem}

\begin{proof}
Theorem~\ref{thm:order-pinning} gives
\(\operatorname{ord}_r(2)/2 = p\) for every prime factor \(r > 3\)
of \(W_p\). The disjointness statement is the contrapositive: if
\(r > 3\) divides both \(W_{p_1}\) and \(W_{p_2}\) then
\(\operatorname{ord}_r(2) = 2 p_1 = 2 p_2\), contradicting
\(p_1 \neq p_2\).
\end{proof}

\begin{remark}
Multi-Factor Pinning is a \emph{global} constraint relating all
prime factors of a single \(W_p\) to a common ``pinned'' order
parameter, and the same pinning applies across different \(W_p\)
to block any shared prime factor.
\end{remark}

\subsection{Phantom Exclusion of the finite danger catalog}

By ``danger triple at \(d\)'' we mean a triple \((r, p, d)\) with
\(r\) an inert prime factor of \(W_p\),
\(d = \operatorname{ord}_r(\omega_3)/4\), and
\((r, p, d)\) surviving every constraint of \S6 (parity, Pell
exclusion, Condition~(II) pass at \(r\)); see
Definition~\ref{def:danger-triple}. A primitive-divisor survey
over all \(d \leq 1000\) admissible under the parity obstruction
of~\S\ref{sec:nct} identifies a finite catalog of such triples
(the ``finite catalog'') with at most five \(d\)-values, including
\(d = 57, 67\), and none with \(d > 683\); see \S\ref{sec:multi-factor}
and \S9.8 for the catalog structure and its secondary-factor
status.

\begin{theorem}[Phantom Exclusion]\label{thm:phantom-exclusion}
Suppose \(W_p\) is composite, every prime factor is inert, and
Condition~(II) holds at every factor (so
Conjecture~\ref{conj:sufficiency} fails at \(W_p\)). Let \(r\)
and \(r'\) be two distinct prime factors of \(W_p\) belonging
respectively to danger triples \((r, p, d)\) and \((r', p', d')\)
in the finite catalog. Then \(p = p'\) and the pair
\((d, d')\) is constrained to a list whose entries share no common
\(p\)-value in the finite catalog.
In particular, the entire finite catalog of danger triples with
\(d \leq 1000\) is \emph{unconditionally} excluded: no composite
\(W_p\) can have two distinct inert factors whose danger triples
both lie in the catalog.
\end{theorem}

\begin{proof}
Both \(r\) and \(r'\) divide \(W_p\), so by
Theorem~\ref{thm:multi-factor-pinning}
\(\operatorname{ord}_r(2)/2 = \operatorname{ord}_{r'}(2)/2 = p\);
in particular \(p = p'\). The primitive-divisor survey of
\S\ref{sec:multi-factor} enumerates the entire catalog of danger
triples \((d, p, r)\) with \(d \leq 1000\) admissible: exactly five
triples, at
\begin{align*}
  (d, p) \in \{\,&(57, 11),\ (57,\, 85\,684\,865\,933),\\
               &(67,\, 24\,822\,855\,876\,938\,710\,967),\\
               &(331, 17),\ (683, 13)\,\}.
\end{align*}
All five \(p\)-values are pairwise distinct. Since a common
\(W_p\) requires a common pinned~\(p\), no two distinct factors
of~\(W_p\) can simultaneously lie in the catalog.
\end{proof}

\begin{remark}[Three-factor reduction + catalog]
Combining Theorem~\ref{thm:multi-factor} (any all-inert composite
has at least three prime factors) with Phantom Exclusion: at least
three distinct factors \(r_1, r_2, r_3\) of \(W_p\) must occur,
all sharing the pinned \(p\)-value, and at most one of these may
belong to the finite catalog. The remaining two are forced to
have \(d > 1000\).
\end{remark}

\subsection{Platinum Lemma for factors below \texorpdfstring{$10^{12}$}{10\textasciicircum 12}}

For computational reasons the ``small \(r\)'' tail is isolated as
a separate unconditional input.

\begin{theorem}[Platinum Lemma]\label{thm:platinum}
No composite \(W_p\) with \(p < 5 \times 10^{11}\) has an inert
prime factor \(r\) with \(r \leq 10^{12}\) and
\(d_r > 43\) that passes Condition~(II).
\end{theorem}

\begin{proof}[Verification sketch]
The statement is a finite computational verification. Each inert
prime factor \(r \leq 10^{12}\) and pinned \(p =
\operatorname{ord}_r(2)/2 < 5 \times 10^{11}\) is enumerated from
the \((r, p)\)-table described in \S9.5. For each \((r, p)\)-row
the Pell-sequence exclusion
(Theorem~\ref{thm:pell-exclusion}) disposes of \(d_r \leq 43\); for
the remaining rows, direct evaluation of the residue
\(\omega_3^{2 W_{p-2}} \pmod r\) either exhibits a
Condition~(II) failure or witnesses a pass, and none of the
\(684\,965\,381\) enumerated rows yields a pass with
\(d_r > 43\). Full scripts and certificates are archived
(see ``Supplementary material'').
\end{proof}

\begin{remark}
The Platinum Lemma is the only computational input in the
cross-case reductions of this section; every downstream
consequence (including
Theorem~\ref{thm:second-moment-reduction}) is unconditional once
the enumeration is accepted. Extensions to \(r \leq 10^{13}\) are
in progress and would narrow the pair-count residual but are not
required for the present reduction.
\end{remark}

\subsection{Exact arithmetic-progression density}

\begin{theorem}[Exact AP density]\label{thm:exact-ap}
For every odd prime \(q > 3\) with \(q \in A\) (equivalently,
\(v_2(\operatorname{ord}_q(2)) = 1\)), set
\(m_q := \operatorname{ord}_q(2)/2\), which is odd. Then for every
prime \(p \geq 5\):
\[
  q \mid W_{p-2} \quad\Longleftrightarrow\quad
  q \in A \text{ and } p \equiv m_q + 2 \pmod{2 m_q}.
\]
In particular, for each \(q \in A\) the set of admissible \(p\)
forms a single arithmetic progression of period \(2 m_q\), and
for \(q \notin A\) the divisibility \(q \mid W_{p-2}\) is
impossible.
\end{theorem}

\begin{proof}
The divisibility \(q \mid W_{p-2}\) is equivalent to
\(q \mid 2^{p-2} + 1\), i.e.\ \(2^{p-2} \equiv -1 \pmod q\). The
latter forces \(\operatorname{ord}_q(2)\) to divide \(2(p-2)\)
without dividing \(p-2\); hence
\(v_2(\operatorname{ord}_q(2)) = 1\) and \(q \in A\). Writing
\(\operatorname{ord}_q(2) = 2 m_q\) with \(m_q\) odd, the element
\(2^{m_q}\) has order exactly \(2\) in
\((\mathbb{Z}/q)^{\times}\) and hence equals \(-1\) (the unique
element of order \(2\) in a cyclic group of even order). So
\(2^{p-2} \equiv -1 \equiv 2^{m_q} \pmod q\) holds iff
\(p - 2 \equiv m_q \pmod{2 m_q}\), i.e.\ \(p \equiv m_q + 2
\pmod{2 m_q}\). Conversely, \(q \in A\) together with this
arithmetic-progression condition directly yields \(q \mid W_{p-2}\).
\end{proof}

Theorem~\ref{thm:exact-ap} is the arithmetic input used below: it
gives an exact, not merely density-based, description of which
primes \(q\) divide \(W_{p-2}\) as a function of \(p\).

\subsection{Second-Moment Reduction}

Write \(E_3\) for the cardinality of
\[
  \bigl\{\,p \geq 5 \text{ prime} : W_p \text{ composite, every prime factor inert,
  Condition~(II) holds at every factor}\,\bigr\},
\]
the set of primes~\(p\) for which Conjecture~\ref{conj:sufficiency}
could fail through the all-inert branch. For each~\(p\) put
\[
  T_p^{>} \;:=\; \{\,r \text{ prime} : r \mid W_p,\
  r \equiv 3 \pmod 8,\ r > 10^{12},\ d_r > 43,\
  d_r \mid W_{p-2}\,\},
\]
the set of \emph{residual inert danger factors} at~\(p\) --- those
factors not already excluded by the Pell exclusion
(Theorem~\ref{thm:pell-exclusion}), by the Platinum Lemma
(Theorem~\ref{thm:platinum}), or by Phantom Exclusion
(Theorem~\ref{thm:phantom-exclusion}). Both \(E_3\) and each
\(|T_p^{>}|\) are deterministic non-negative integers
(with \(E_3\) possibly \(+\infty\)).

\begin{theorem}[Second-Moment Reduction]\label{thm:second-moment-reduction}
Unconditionally,
\[
  E_3 \;\leq\; \tfrac12 \sum_{p \geq 5} |T_p^{>}|\,(|T_p^{>}| - 1)
           \;=\; \sum_{p \geq 5} \binom{|T_p^{>}|}{2}
           \;=\; \Pi,
\]
where \(\Pi\) is the total number of unordered pairs
\(\{r_1, r_2\}\) of distinct residual inert danger factors of a
single composite Wagstaff number.
\end{theorem}

\begin{proof}
\emph{Step 1 (lower bound on \(|T_p^{>}|\) when \(p\) contributes to~\(E_3\)).}
Let \(p\) be counted by~\(E_3\). By
Theorem~\ref{thm:multi-factor}, the all-inert composite \(W_p\) has
at least three distinct prime factors \(r_1, r_2, r_3 > 3\), each
inert with \(r_i \equiv 3 \pmod 8\) and each a danger factor
(i.e.\ \(d_{r_i} \mid W_{p-2}\) and Condition~(II) holds locally
at \(r_i\)). The Pell exclusion
(Theorem~\ref{thm:pell-exclusion}) forces \(d_{r_i} > 43\) for
each~\(r_i\), and Phantom Exclusion
(Theorem~\ref{thm:phantom-exclusion}) rules out two of the factors
simultaneously belonging to the finite catalog at
\(d \leq 1000\); hence at most one of \(r_1, r_2, r_3\) has
\(d_{r_i} \leq 1000\). For the remaining two factors the Platinum
Lemma (Theorem~\ref{thm:platinum}) gives \(r_i > 10^{12}\), so
both lie in~\(T_p^{>}\). Therefore \(|T_p^{>}| \geq 2\).

\emph{Step 2 (elementary integer inequality).}
For every non-negative integer~\(X\),
\[
  \mathbf{1}_{X \geq 2} \;\leq\; \binom{X}{2} \;=\; \tfrac{X(X-1)}{2},
\]
as is verified directly for \(X \in \{0, 1, 2, 3, 4\}\) and holds a
fortiori for \(X \geq 3\).

\emph{Step 3 (summation).}
By Step~1, \(E_3 \leq \#\{\,p : |T_p^{>}| \geq 2\,\} = \sum_p
\mathbf{1}_{|T_p^{>}| \geq 2}\). Applying Step~2 termwise,
\[
  E_3 \;\leq\; \sum_{p \geq 5} \mathbf{1}_{|T_p^{>}| \geq 2}
     \;\leq\; \sum_{p \geq 5} \binom{|T_p^{>}|}{2}
     \;=\; \Pi. \qedhere
\]
\end{proof}

\begin{corollary}[Residual pair-count bound]\label{cor:residual-pair-count}
If \(\Pi < 1\) --- equivalently, if no composite Wagstaff number
\(W_p\) admits two distinct residual inert danger factors \(r_1,
r_2 > 10^{12}\) at the same~\(p\) --- then \(E_3 = 0\), and
Conjecture~\ref{conj:sufficiency} holds unconditionally in the
all-inert branch. More generally, any finite bound on~\(\Pi\)
yields a finite bound on~\(E_3\).
\end{corollary}

\begin{remark}[No independence or heuristic hypothesis]
Every step above is deterministic: neither the theorem nor its
corollary invokes any probability measure, independence, or
randomness heuristic on prime factors. The step from a
three-factor obstruction to a pair-count obstruction uses only the
integer inequality \(\mathbf{1}_{X\geq 2} \leq \binom{X}{2}\), and
the conversion to the pair count~\(\Pi\) is the definition of
\(\binom{X}{2}\). This is in contrast with the heuristic
third-moment counts employed in earlier checkpoints, which are
inputs to the conjectural density estimate on~\(\Pi\) discussed
in~\S10, not to Theorem~\ref{thm:second-moment-reduction} itself.
\end{remark}

\begin{remark}[Structure of \(T_p^{>}\) and shared pinning]
By Multi-Factor Pinning
(Theorem~\ref{thm:multi-factor-pinning}), every element of~\(T_p^{>}\)
satisfies \(\operatorname{ord}_r(2) = 2p\) for the \emph{same}
value of~\(p\), so a pair \(\{r_1, r_2\} \subset T_p^{>}\) is
more tightly constrained than a pair of independent danger primes:
both factors share the pinned order parameter and the exact AP
(Theorem~\ref{thm:exact-ap}) admissibility congruence class on~\(p\)
in terms of \(d_{r_1}\) and~\(d_{r_2}\). This is the shape of the
residual problem in~\S10.
\end{remark}

\subsection{The residual problem}

The chain of reductions developed above transfers
Conjecture~\ref{conj:sufficiency}, in its all-inert branch, to the
single deterministic claim
\begin{equation}\label{eq:residual-target}
  \Pi \;:=\; \sum_{p \geq 5} \binom{|T_p^{>}|}{2} \;=\; 0,
\end{equation}
that is, no composite Wagstaff number admits a pair of distinct
residual inert danger factors. This is a pair-counting question on
a thin sieve set of inert primes; see \S10 for candidate approaches
and the structural obstacles identified in prior work.

\begin{proposition}[Residual problem status]\label{prop:residual-status}
Every composite \(W_p\) with \(p < 5 \times 10^{11}\) and every
Wagstaff probable prime with \(p < 100\,000\) tested to date is
consistent with \eqref{eq:residual-target}: each has
\(|T_p^{>}| = 0\) (no residual factor at all) or
\(|T_p^{>}| = 1\) (a single residual factor, disposed of by the
split/secondary-factor analysis of~\S6.7 and~\S9.8).
\end{proposition}

\begin{proof}[Verification]
The \(684\,965\,381\)-row Platinum enumeration of \S9.5 records
every inert factor below \(10^{12}\); combined with the
Platinum Lemma and the \(d \leq 43\) exclusion, the empirical
count \(|T_p^{>}|\) on the enumerated range is \(0\) or \(1\) at every
\(p\) tested, and no pair \(\{r_1, r_2\} \subset T_p^{>}\) has
been found in the enumerated range.
\end{proof}

\section{Computational verification}


The software toolchain used in this section (SymPy, gmpy2,
PARI/GP) is described in \S9.7.

\subsection{Proven Wagstaff primes}


Among the 36 known Wagstaff primes and probable primes, those with
\(p \leq 42737\) are proved prime --- the smaller exponents by
classical methods, the larger ones (\(p \geq 2617\)) by the
elliptic curve primality proving (ECPP) algorithm of Atkin--Morain
\cite{ref1}. The remaining exponents (\(p \geq 83339\)) are
probable primes only. This paper provides proofs by different
methods for a subset:

\begin{center}
\begin{tabular}{lll}
\hline
Range & Method in this paper & Notes \\
\hline
\(p = 3\) & Trivial &  \\
\(5 \leq p \leq 1709\) (20 primes) & \(N+1\) criterion (Thm~\ref{thm:criterion}) &  \\
\(p = 2617\) & BLS \(N-1\) (Thm~\ref{thm:W2617}) & Also proved by ECPP \\
\(p = 10501\) & BLS \(N-1\) (Thm~\ref{thm:W10501}) & Also proved by ECPP \\
\(p = 12391\) & BLS \(N-1\) (Thm~\ref{thm:W12391}) & Also proved by ECPP \\
\hline
\end{tabular}
\end{center}


The BLS proofs for \(W_{2617}\), \(W_{10501}\), and \(W_{12391}\)
are independent of ECPP: each reduces to a list of BLS witness
pairs together with APR-CL primality certificates for the
factors, all verifiable by a short computation.
All 36 known Wagstaff primes and probable primes satisfy Condition (II).

\subsection{Conjecture~\ref{conj:R3} verification}


All 185 known primes \(r \equiv 1 \pmod{8}\) with \(r \mid W_p\)
satisfy Conjecture~\ref{conj:R3}. The distribution by proof method:

\begin{center}
\begin{tabular}{lll}
\hline
Configuration & Fraction & Reference \\
\hline
Case A (\(r \mid U_p\)) & 1\% (1) & Theorem~\ref{thm:case-A} \\
Case B (\(r \mid V_p\)) & 3\% (6) & Theorem~\ref{thm:case-B} \\
Case C, \(4U_p^2 \equiv 1\) & 25\% (46) & Theorem~\ref{thm:case-4U} \\
Case C, \(d\) inadmissible & 23\% (43) & Prop~\ref{prop:rank-equivalence} \\
Case C, general (\(d > 200\)) & 48\% (89) & R3 verified \\
\hline
\end{tabular}
\end{center}


Exactly one Case A instance occurs (\(p = 3361\), \(r = 53777\)).
For the 43 ``\(d\) inadmissible'' factors, either
\(4 \nmid \operatorname{ord}_r(\omega_3)\) or
\(\operatorname{ord}_r(\omega_3)/4\) is even; in either case
Proposition~\ref{prop:rank-equivalence} shows Condition~(II) cannot hold.
The remaining 89 Case C factors all have odd
\(d = \operatorname{ord}_r(\omega_3)/4 > 200\); for these,
Conjecture~\ref{conj:R3} is verified by direct computation of
\(\operatorname{ord}_r(1+s)\).

\subsection{Primitive divisor survey}


A parallel computation searched all primes \(q \leq 1000\) with
\(q \equiv 3 \pmod{8}\) and primitive divisors \(r < 10^{12}\) of
\(U_{4q}\) with \(r \equiv 1 \pmod{8}\). Among 62 primitive divisors,
exactly one is a \(q\)-th power case (\(q = 43\), \(r = 14449\)).
An extended ECM computation found
\(r = 12036824328615957881\) for \(q = 163\): 2 is a 163rd power
residue mod \(r\), yet \(v_2(\operatorname{ord}_r(2)) = 1\), so
the natural strengthening in Remark~\ref{rem:qpro-false} fails. Since
\(\operatorname{ord}_r(2)/2\) is composite, NCT is unaffected.

\subsection{NCT verification}


No compatible triple \((d, r, p)\) exists for any odd
\(d \leq 81\); this is the unconditional Theorem~\ref{thm:NCT-81}. For the
parity-unblocked range \(83 \leq d \leq 200\) --- the 11 values
\(d \in \{83, 99, 107, 121, 129, 131,\)
\(139, 163, 171, 177, 179\}\)
surviving the parity obstruction of Theorem~\ref{thm:NCT-81} --- an independent
sieve-based verification was carried out on a 128-core machine.
For each such \(d\), all primes \(r = 8dk + 1\) with \(r < 10^{11}\)
were enumerated, the rank-of-apparition condition
\(\rho(r) = 4d\) was tested via the Pell recurrence modulo \(r\),
and surviving primitive divisors were checked for compatibility with
some prime \(p\) satisfying \(\operatorname{ord}_r(2) = 2p\) and
\(d \mid W_{p-2}\). The enumeration tested
\(\sim 1.07 \times 10^8\) primes, found 12 primitive Pell divisors
(at \(d \in \{121, 131, 163, 171, 177, 179\}\)), and produced
zero compatible triples. For \(d > 200\), complete factorization of
\(U_{4d}\) becomes the computational bottleneck; the sieve method
above scales further but at increasing cost.

\subsection{Inert factor survey}


\emph{Local versus global pass.} We say a composite \(W_p\)
admits a \emph{local pass} of Condition (II) at a prime divisor
\(r \mid W_p\) if \(\omega_3^{(W_p+1)/2} \equiv -1 \pmod r\); a
\emph{global pass} means Condition (II) holds at every prime
divisor of \(W_p\). Conjecture~\ref{conj:sufficiency} is the statement that no
composite \(W_p\) admits a global pass. A single local pass does
not by itself produce a counterexample: it only asserts that the
composite \(W_p\) survives the local test at one prime divisor,
with the remaining cofactor untested.

Among 869 known prime factors \(r \equiv 3 \pmod{8}\) of composite
\(W_p\) with \(r < 2 \times 10^5\): Condition (II) fails at every
factor.

\begin{center}
\begin{tabular}{llll}
\hline
\(G_r\) range & Count & Fraction & Status \\
\hline
\(G_r = 1\) & 235 & 27.0\% & Corollary~\ref{cor:small-Gr} \\
\(G_r \leq 3\) & 829 & 95.4\% & Corollary~\ref{cor:small-Gr} \\
\(G_r \leq 43\) & 845 & 97.2\% & Corollary~\ref{cor:small-Gr} \\
\(G_r > 43\) & 24 & 2.8\% & Verified \\
\hline
\end{tabular}
\end{center}


For the 24 factors with \(G_r > 43\):
\(\operatorname{ord}_r(\omega_3)/4\) exceeds \(G_r\) in every case,
with ratio \(\operatorname{ord}/4 \div G_r \geq 5\). The excess
arises because \(\operatorname{ord}_r(\omega_3)/4\) contains a prime
factor \(q \notin A\) (equivalently, \(q \nmid W_{p-2}\)), preventing
\(\operatorname{ord}/4\) from dividing \(W_{p-2}\).

The 11 distinct \(G_r\) values observed are
\(\{1, 3, 9, 11, 19, 57, 129, 177, 321, 2217, 2451\}\); in particular,
all prime factors of observed \(G_r\) lie in
\(\{3, 11, 19, 43, 59, 107, 739\}\), each of which is Class III,
consistent with the structure of \(A\).

\textbf{Extended inert-factor (Platinum) survey.} A segmented
sieve enumerated every prime \(r \equiv 3 \pmod{8}\) with
\(r < 10^{12}\), computed \(\operatorname{ord}_r(2)\), and
recorded every pair \((p,r)\) with
\(\operatorname{ord}_r(2) = 2p\) (so \(r \mid W_p\)) and \(W_p\)
composite. The enumeration produced 684{,}965{,}381 such pairs,
with \(p < 5 \times 10^{11}\). Condition (II) was tested at each
factor. Of these, 97.72\% (669{,}378{,}360 pairs) satisfy
\(G_r \leq 43\) and are excluded by Corollary~\ref{cor:small-Gr}; the remaining
2.28\% (15{,}587{,}021) are excluded individually by the
order-excess argument of \S6.5. No composite \(W_p\) with all
factors passing Condition (II) was observed. (Condition (II)
necessarily passes at factors of Wagstaff primes, by
Proposition~\ref{prop:base-3}; such pairs are excluded from the count.)
This enumeration is the computational foundation of
Theorem~\ref{thm:platinum} (Platinum Lemma) and of
Proposition~\ref{prop:residual-status}: combined with the
structural exclusions of \S\ref{sec:pinning-reductions}, no inert
factor with \(r \leq 10^{12}\) contributes to an open
counterexample configuration.
Operational details --- parallelisation, segmented sieve scale,
the three local-pass witnesses, and the secondary-factor closures
of \S10 Problem 2 --- are recorded in \S9.8.

\subsection{The Chebyshev test algorithm}


\emph{Unconditionally, the algorithm below is a probable-prime test for
the family \(W_p\)}: it outputs \textbf{probable prime} for every
Wagstaff prime (Proposition~\ref{prop:base-3}) and has never been observed to
output \textbf{probable prime} on a composite \(W_p\) (\S9.1, \S9.5).
Condition (I) is automatically satisfied (Theorem~\ref{thm:cond-I}) and need
not be checked. Conjecture~\ref{conj:sufficiency} would upgrade the test to a
deterministic standalone primality test; the unconditional results
of this paper (Theorems~\ref{thm:cond-I}, \ref{thm:NCT-81}, \ref{thm:pell-exclusion} and Lemmas~\ref{lem:mod-8-exclusion}, \ref{lem:v2-classification}) reduce
that conjecture to the two structurally independent gaps stated in
\S10.

\begin{quote}
\textbf{Algorithm} (Wagstaff Chebyshev test).
\textbf{Input:} An odd prime \(p \geq 5\).
1. Set \(N \leftarrow (2^p + 1)/3\) and
\(\omega \leftarrow 3 + 2\sqrt{2}\) in \(\mathbb{Z}[\sqrt{2}]/(N)\).
2. Compute \(R \leftarrow \omega^{(N+1)/2} \bmod N\) via binary
exponentiation in \(\mathbb{Z}[\sqrt{2}]/(N)\).
3. If \(R = -1\): output \textbf{probable prime}.
Else: output \textbf{composite}.
\end{quote}

The test requires one modular exponentiation in
\(\mathbb{Z}[\sqrt{2}]/(N)\), i.e. \(O(p)\) modular squarings of
\(O(p)\)-bit integers, for a total cost of \(\widetilde{O}(p^2)\)
bit operations using FFT-based integer multiplication; with
schoolbook multiplication this becomes \(O(p^3)\). The constant in
\(\widetilde{O}\) is the same as for one modular exponentiation
\(\bmod\, N\) of an integer the size of \(N\).

\subsection{Computer-assisted components}


Several theorems in this paper have proofs that combine algebraic
arguments with finite computations. We list each such result, the
exact computation required, and its scope.

\begin{center}
\begin{tabular}{lp{6.2cm}p{4.5cm}}
\hline
Result & Computation required & Scope \\
\hline
Theorem~\ref{thm:criterion} & Factorization of \(W_{p-2}\) & Each application (\S9.1) \\
Theorem~\ref{thm:W2617} & BLS certificate for \(W_{2617}\) & Single number \\
Theorem~\ref{thm:W10501} & BLS certificate for \(W_{10501}\); APR-CL primality of 103 factors & Single number \\
Theorem~\ref{thm:W12391} & BLS certificate for \(W_{12391}\); APR-CL primality of 61 factors & Single number \\
Theorem~\ref{thm:NCT-81} & Complete factorization of \(U_{4d}\) for each admissible \(d \leq 81\); for each primitive \(r \equiv 1\,(8)\), verify \(v_2(\operatorname{ord}_r(2)) \neq 1\) or \(\operatorname{ord}_r(2)/2\) composite or \(d \nmid W_{p-2}\) & 10 values of \(d\); largest factorization: \(U_{324}\) (124 digits) \\
Theorem~\ref{thm:pell-exclusion} & Complete factorization of \(N_d = V_d^2+2\) for each admissible \(d \leq 43\); for each prime \(r \equiv 3\,(8)\) dividing \(N_d\), verify the same three conditions & 6 values of \(d\); largest factorization: \(N_{43}\) (33 digits) \\
Theorem~\ref{thm:multi-factor} & Complete factorization of \(V_{114}\), \(V_{134}\); partial factorization of \(V_{662}\); primitive-divisor survey of all 79 admissible \(d \leq 1000\) & 4 values of \(d\) with danger triples \\
\hline
\end{tabular}
\end{center}


The algebraic arguments (parity obstruction, Theorems~\ref{thm:case-A}--\ref{thm:case-4U},
\ref{thm:quartic-residuosity}, Corollary~\ref{cor:small-Gr}, Propositions~\ref{prop:rank-equivalence} and~\ref{prop:qth-power}, Theorem~\ref{thm:generalized-LTE}, and
the results of Appendix A) are fully unconditional and require no
computation. The finite verifications above are deterministic and
reproducible; factorization certificates and verification logs are
in the supplementary repository.

\textbf{Software toolchain.} All computational components in this paper
are reproducible from the scripts in the supplementary repository
using free, open-source software.

\begin{itemize}
\item \emph{Integer factorization} (cyclotomic factors \(\Phi_d(2)\) for the
  BLS proofs, the Pell numbers \(U_{4d}\), \(V_{4d}\), and the Pell
  exclusion numbers \(N_d\)): SymPy 1.12+ \texttt{factorint} (combining trial
  division, Pollard \(\rho\), Pollard \(p\!-\!1\), and ECM internally),
  augmented for cyclotomic factors by trial division along the
  arithmetic progression \(r \equiv 1 \pmod{d}\) implied by primitive
  divisor theory.
\item \emph{Pre-known large factorizations}: the Cunningham project tables
  \cite{ref4} for cyclotomic
  factors \(\Phi_d(2)\) with \(d \leq 1200\), and FactorDB as a
  discovery aid for two larger terms
  (\(\Phi_{2065}(2)\) and \(\Phi_{2478}(2)\) in
  Theorem~\ref{thm:W12391}). Every retrieved factor is independently
  re-verified for primality by APR-CL before use; FactorDB plays no
  evidentiary role.
\item \emph{Primality of cofactors} (all factors used in BLS certificates):
  the APR-CL implementation in PARI/GP 2.15+
  (\texttt{isprime(n)}), which is unconditionally deterministic and
  follows the Adleman--Pomerance--Rumely / Cohen--Lenstra algorithm
  \cite{ref6}. The largest cofactor is a 371-digit factor of
  \(\Phi_{2065}(2)\) (Theorem~\ref{thm:W12391}); the smallest are single-digit
  primes, but all receive APR-CL certification for uniformity.
\item \emph{Fast modular arithmetic} (the modular exponentiations
  \(a^{N-1} \bmod N\) and \(\omega_3^{(N+1)/2} \bmod N\) for
  \(N = W_{10501}, W_{12391}\)): the GMP-backed bindings in
  gmpy2 2.1+, falling back to SymPy when gmpy2 is unavailable.
\end{itemize}

The driver scripts
\texttt{bls\_n\_minus\_1\_w\{2617,10501,12391\}.py}
in the supplementary repository emit JSON certificates
(\texttt{bls\_certificate\_w*.json}) recording every cyclotomic factor,
prime decomposition, BLS witness, and primality method used,
along with full execution logs.

\subsection{Survey scale, witnesses, and secondary-factor closures}


The extended inert-factor survey of \S9.5 was carried out on a
128-core machine. The segmented sieve scanned
\(1.93 \times 10^{10}\) primes in \(\approx 3.98\) hours of
wall-clock time, producing the 684{,}965{,}381 pairs \((p,r)\)
described in \S9.5.
We summarise four outputs: the local-pass witnesses (\S9.5
found three), the companion split sweep, and the secondary-factor
closures that dispose of the two remaining gap configurations.

\emph{Three local-pass witnesses.} Of the 684{,}965{,}381 inert
pairs, exactly three survive all single-factor local
obstructions --- i.e.\ exhibit a local pass of Condition (II) at
\(r\). In each case the pass is only local: the remaining
cofactor of \(W_p\) is unfactored and untested. The three
witnesses cover structurally distinct mechanisms:

\begin{enumerate}
\item \emph{Primitive-root coincidence.}
  \(p = 1{,}251{,}488{,}009\), \(r = 2p+1\),
  \(d = G_r = 41{,}716{,}267\). Here \(r\) is the smallest prime
  in the progression \(r \equiv 1 \pmod{2p}\); the coincidence
  has no further algebraic content.
\item \emph{Structured cofactor.}
  \(p = 10{,}916{,}765{,}939\), \(r = 152{,}834{,}723{,}147\),
  \(d = G_r = 38{,}208{,}680{,}787\). All three prime factors of
  \(d\) (namely \(3, 5419, 2350291\)) lie in the Class III set
  \(A\), so the order parameter is accommodated inside the
  multiplicative structure of \(W_{p-2}\).
\item \emph{Small-\(d\) gap at \(d = 57\).}
  \(p = 85{,}684{,}865{,}933\), \(r = 2p+1\),
  \(d = G_r = 57 = 3 \cdot 19\). This is the single-factor
  residual configuration at \(d = 57\) identified in \S10,
  Problem 2; it is disposed of globally by the secondary-factor
  closure below.
\end{enumerate}

For witnesses 1 and~2, \(W_p\) has \(\approx 10^{8.6}\) and
\(\approx 10^{9.5}\) digits respectively; no secondary factor is
known, and their global disposition is genuinely open --- they
illustrate the structural gap at large~\(d\) discussed in~\S10,
Problem~2.

\emph{Companion split sweep.}
Full factor data for each witness and a companion Conjecture~\ref{conj:R3} check on
all \(14{,}812{,}084\) split prime divisors with
\(r \leq 10^{11}\) (zero failures\footnote{The companion sweep
scanned \(4.12 \times 10^9\) primes \(r \equiv 1 \pmod 8\) and
verified \(p \mid \operatorname{ord}_r(\alpha)\) on every split
pair; under the null in which \(\chi_r(\alpha)\) projects
uniformly onto the \(p\)-th roots of unity, the expected failure
count is \(E = \sum 1/p \approx 0.76\) and
\(\exp(-E) \approx 0.47\). The observation therefore adds no
independent evidence for Conjecture~\ref{conj:R3} beyond the unconditional closures of
Theorem~\ref{thm:NCT-81}; we report it only for
completeness.}) are recorded in the Zenodo archive.

\emph{Secondary-factor closures (\S10, Problem 2).} The dangerous
single-factor configurations at \(d = 57\) and \(d = 67\) are
ruled out at secondary prime factors of the same composite \(W_p\):

\begin{itemize}
\item For \(d = 57\): the prime
  \(r_2 = 7{,}675{,}646{,}348{,}774{,}307{,}083 = 2pk + 1\)
  with \(k = 44{,}789{,}977\) divides \(W_{85684865933}\),
  is \(\equiv 3 \pmod{8}\) (inert), and has
  \(\operatorname{ord}_{r_2}(\omega_3) = r_2 + 1\). The integer
  \((r_2+1)/4\) has the prime divisor \(487\), which lies outside
  \(A\) since \(\operatorname{ord}_{487}(2) = 243\) is odd. By
  Theorem~\ref{thm:parity-obstruction}, Condition (II) fails at \(r_2\), and hence globally
  on \(W_p\).
\item For \(d = 67\): the prime
  \(r_2 = 5{,}474{,}333{,}343{,}676{,}555{,}561{,}818{,}313 = 2pk+1\)
  with \(k = 110268\) divides \(W_{24822855876938710967}\), and
  is \(\equiv 1 \pmod{8}\) (split). Conjecture~\ref{conj:R3} holds at \(r_2\)
  (verified directly by computing
  \(\omega_3^{(r_2-1)/p}\) modulo \(r_2\)), so by Theorem~\ref{thm:R3-conditional}
  Condition (II) fails at \(r_2\) and hence globally on \(W_p\).
\end{itemize}

These two configurations are disposed of computationally at
secondary factors. They are the only two single-factor
configurations with \(d \leq 1000\) that the single-factor
argument of \S\S5--6 currently fails to rule out, and each is
dispositioned at a second factor of the same composite \(W_p\).
They are not counted among the unconditional results in the
status table of \S1.2.
Theorem~\ref{thm:phantom-exclusion} further excludes any two
such catalog triples from coexisting in the same composite
\(W_p\); together with Theorem~\ref{thm:multi-factor}, a
counterexample \(W_p\) would therefore require at least two inert
factors outside the \(d \leq 1000\) catalog and the residual
pair-count bound of Theorem~\ref{thm:second-moment-reduction}.

\emph{Danger-triple enumeration (\S\ref{sec:multi-factor}).} The complete factorisation
of \(V_{114}\) and \(V_{134}\) and the partial factorisation of
\(V_{662}\), together with a primitive-divisor survey of all 79
admissible \(d \leq 1000\), identify five danger triples
(Definition~\ref{def:danger-triple}): the two real configurations
above, and three phantom triples at \((d, p) = (57, 11)\),
\((331, 17)\), and \((683, 13)\), where in each case \(W_p\) is
prime. No other admissible \(d \leq 1000\) yields a danger triple
among the identified primitive divisors.
The factorisation data are in the Zenodo archive.

\begin{remark}[Single-factor versus global sufficiency]\label{rem:single-vs-global}
In contrast to the Lucas--Lehmer test for Mersenne numbers, where the
congruence \(s_{p-2} \equiv 0 \pmod{M_p}\) fails independently at
each prime factor of a composite \(M_p\), the Chebyshev criterion
for Wagstaff numbers admits local passes of Condition~(II) at
individual factors (witnesses 1--3 above). Global failure is then
recovered only through cross-factor arguments at secondary
divisors of the same \(W_p\). A complete sufficiency proof
therefore cannot proceed factor by factor; it requires a global
argument that accounts for the interaction between distinct prime
divisors of a composite \(W_p\).
Theorem~\ref{thm:multi-factor} gives one such global argument: the
semiprime exclusion and the finiteness of danger triples reduce the
residual gap from ``any inert factor with \(d > 43\) might pass'' to
``at least one factor must have a danger triple at some
\(d \notin \{57, 67, 331, 683\}\),'' with no such triple identified
for \(d \leq 1000\).
\end{remark}



\section{Open problems}

\textbf{Problem 1} (NCT for general \(d\)). Prove that no compatible triple
\((d, r, p)\) exists for any odd \(d > 1\). Proved for \(d \leq 81\)
(Theorem~\ref{thm:NCT-81}); verified computationally through \(d = 200\) (\S9.4).
The parity obstruction eliminates 80\% of \(d\); the remaining
cases require a direct compatible-triple argument (Proposition~\ref{prop:qth-power},
Remark~\ref{rem:qpro-false}). See Remark~\ref{rem:split-gap} for the structural source of this gap.

\textbf{Problem 2} (Residual pair-count bound at inert factors). For
\(r \equiv 3 \pmod{8}\) dividing a composite \(W_p\), prove that
\(\operatorname{ord}_r(\omega_3)/4 \nmid W_{p-2}\). The unconditional
reductions of \S\ref{sec:pinning-reductions} transfer this branch
to a single deterministic residual. Specifically:
Theorem~\ref{thm:pell-exclusion} handles all admissible
\(d \leq 43\); Theorem~\ref{thm:phantom-exclusion} (Phantom
Exclusion) unconditionally rules out the entire finite catalog of
danger triples at admissible \(d \leq 1000\) produced by the
primitive-divisor survey of \S\ref{sec:multi-factor}; and
Theorem~\ref{thm:platinum} (Platinum Lemma) excludes every factor
\(r \leq 10^{12}\). After these reductions the only remaining
input is the Second-Moment Reduction target
(Theorem~\ref{thm:second-moment-reduction} and
Corollary~\ref{cor:residual-pair-count}):
\[
  \Pi \;:=\; \sum_{p \geq 5} \binom{|T_p^{>}|}{2} \;=\; 0,
\]
i.e.\ no composite Wagstaff number has two distinct inert prime
factors \(r_1, r_2\) with \(r_i \equiv 3 \pmod 8\),
\(r_i > 10^{12}\), \(d_{r_i} > 43\), and
\(d_{r_i} \mid W_{p-2}\). Equivalently, the global unordered
pair-count of residual inert danger factors over all Wagstaff
numbers vanishes.

A heuristic count supports \(\Pi = 0\). Modelling each primitive
inert divisor of \(V_{2d}\) as a random coincidence with
\(\operatorname{ord}_r(2)/2 = p\) among its admissible residues,
the expected total number of \emph{single} danger triples
\((d, r, p)\) with \(d > D\) is
\[
  E(D) \;\approx\; \sum_{\substack{d > D \\ d \text{ admissible}}}
  \frac{(\ln\ln|V_{2d}|)^2}{2\,\ln|V_{2d}| \cdot M_d},
\]
where \(M_d = \operatorname{lcm}(\operatorname{ord}_q(2)/2 : q \mid d)\)
and \(|V_{2d}| \sim (1+\sqrt{2})^{2d}\). The series converges:
for admissible prime \(d = q\), \(M_d = \operatorname{ord}_q(2)/2
\sim q/2\), so the \(d\)-th term is
\(O((\ln\ln d)^2/d^2)\); for composite admissible \(d\) the lcm
\(M_d\) grows at least as fast. Numerical evaluation at
\(D = 1000\) yields \(E \approx 0.003\), consistent with the zero
count observed. For \emph{pair} counts the same heuristic, together
with the Multi-Factor Pinning constraint that the two factors
share the order parameter exactly, gives a total pair count
heuristic of \(O(10^{-6})\), consistent with \(\Pi = 0\); the
rigorous proof of \(\Pi = 0\) is open.

The closed form \(N_d = V_d^2 + 2\)
(Proposition~\ref{prop:Nd-closed-form}) exhibits \(N_d\) as a norm
in \(\mathbb{Z}[\sqrt{-2}]\); whether the Zsigmondy theory of
primitive divisors of such norms can yield an unconditional
pair-count bound is an open question.

\begin{remark}\label{rem:generalization}
The structural ingredients used in \S\S5--6 --- the
\(v_2\) classification of \(\operatorname{ord}_r(\omega_3)\)
(Lemma~\ref{lem:v2-classification}), the Pell exclusion numbers \(N_d = V_d^2+2\) (Proposition~\ref{prop:Nd-closed-form}), and the compatible-triple formulation of Conjecture~\ref{conj:R3}
(Proposition~\ref{prop:rank-equivalence}) --- are properties of the fixed element
\(\omega_3 \in \mathbb{Z}[\sqrt{2}]\) and the Pell sequences. The
Wagstaff-specific content enters through Lemma~\ref{lem:ord-2}
(\(\operatorname{ord}_r(2) = 2p\)) and the identity
\(N+1 = 4W_{p-2}\), both of which have analogues in the
generalised Chua family \(N_b = (b^p+1)/(b+1)\): for each base
\(b\), the Chebyshev element \(\omega_b = b + \sqrt{b^2-1}\) plays
the role of \(\omega_3\), and the identity \(N_b + 1\) factors
similarly. The extension requires \(D = b^2 - 1\) to be non-square,
so the Pell sequences are replaced by the Lucas sequences of
\(\mathbb{Z}[\sqrt{D}]\).
\end{remark}



\appendix
\section{Independent local algebra at split Wagstaff factors}


This appendix gives the first closed-form restriction on the
quadratic-residue behaviour of \(\alpha = 1+\sqrt 2\) at split
Wagstaff factors: a dichotomy determined by \(r \bmod 16\)
(Theorem~\ref{thm:A.8}). The proof assembles a classical quartic-residue
identity in the Pell representation \(r = a^2 - 2b^2\) (\S A.1,
yielding Corollary~\ref{cor:A.5}), a Wagstaff-specific Octic Dichotomy
(Theorem~\ref{thm:A.6}), and a genus formula for
\((1+\sqrt 2 \mid r)\) (Theorem~\ref{thm:A.7}). The appendix is
independent of the main reduction of \S5 --- Conjecture~\ref{conj:R3} at split factors
is handled there by the Pell-rank reformulation
(Proposition~\ref{prop:rank-equivalence}) and Theorem~\ref{thm:NCT-81} --- but provides
additional local information: the QNR
Dichotomy (Theorem~\ref{thm:A.8}) shows that \(\alpha = 1+\sqrt{2}\)
is a quadratic non-residue at every split Wagstaff factor with
\(v_2(r-1) = 3\). Whether this quadratic-character constraint
interacts with the compatible-triple conditions of \S5.6
remains open.

\subsection{The quartic character of 2 in the Pell representation}


Let \(r \equiv 1 \pmod 8\) be a fixed prime and \(r = a^2 - 2b^2\)
a representation with \(a > 0\). Since \(r\) is odd, \(a\) is odd,
and \(\gcd(a,b) = 1\) (any common prime divisor of \(a\) and \(b\)
would also divide \(r\), forcing \(r\) itself to divide \(a\) or
\(b\), absurd since \(0 < a, |b| < r\)). Also \(b\) is even: from
\(a\) odd we have \(a^2 \equiv 1 \pmod 8\), and \(r \equiv 1 \pmod 8\)
forces \(2b^2 \equiv 0 \pmod 8\), hence \(b^2 \equiv 0 \pmod 4\).

\begin{lemma}[Free QR lemma]\label{lem:A.1}
\((b/r) = +1\).
\end{lemma}


\begin{proof}
Write \(b = 2^k b'\) with \(b'\) odd, \(k \geq 1\). Then
\((b/r) = (2/r)^k (b'/r) = (b'/r)\), since \((2/r) = +1\) (because
\(r \equiv 1 \pmod 8\)). Apply quadratic reciprocity to the odd
positive integer \(b'\) and the odd prime \(r\):
\[
(b'/r) \cdot (r/b') = (-1)^{((b'-1)/2)\cdot((r-1)/2)}.
\]
Since \(r \equiv 1 \pmod 4\), \((r-1)/2\) is even, so the right
side is \(+1\) and \((b'/r) = (r/b')\). Now
\(2b^2 = 2^{2k+1} b'^2\) is divisible by \(b'\), so
\(r = a^2 - 2b^2 \equiv a^2 \pmod{b'}\), and
\((r/b') = (a^2/b') = +1\) (using \(\gcd(a,b') = 1\)). Therefore
\((b/r) = +1\).\end{proof}


\begin{lemma}[Euler-criterion lift]\label{lem:A.2}
\((a/r) = 2^{(r-1)/4} \cdot (b/r)\)
in \(\mathbb F_r\).
\end{lemma}


\begin{proof}
From \(r = a^2 - 2b^2\) we have \(a^2 \equiv 2b^2 \pmod r\).
Raise both sides to the \((r-1)/4\) power:
\[
a^{(r-1)/2} \equiv 2^{(r-1)/4} \cdot b^{(r-1)/2} \pmod r.
\]
By Euler's criterion, the two outer exponents reduce to \((a/r)\)
and \((b/r)\) in \(\{\pm 1\}\); the middle factor
\(2^{(r-1)/4}\) is the square root of \((2/r) = +1\), hence also
in \(\{\pm 1\}\), so the congruence is an equality of integers.\end{proof}


\begin{lemma}[Reciprocity flip]\label{lem:A.3}
\((a/r) = (-2/a)\).
\end{lemma}


\begin{proof}
Jacobi reciprocity on the coprime odd positive integers
\(a\) and \(r\) gives \((a/r)(r/a) = (-1)^{((a-1)/2)((r-1)/2)} = +1\),
since \((r-1)/2\) is even. So \((a/r) = (r/a)\). From
\(r = a^2 - 2b^2\) we get \(r \equiv -2b^2 \pmod a\); since
\(\gcd(a,b) = 1\),
\((r/a) = (-2b^2/a) = (-2/a)(b/a)^2 = (-2/a)\).\end{proof}


\begin{proposition}[Pell quartic-character identity, Dirichlet]\label{prop:A.4}
Let \(r \equiv 1 \pmod 8\) be a prime and \(r = a^2 - 2b^2\) any
representation with \(a > 0\). Then in \(\mathbb F_r\),
\[
2^{(r-1)/4} \equiv (-2/a) \pmod r.
\]
Equivalently, 2 is a fourth-power residue modulo \(r\) iff
\(a \equiv 1\) or \(3 \pmod 8\).
\end{proposition}


\begin{proof}
Combine Lemmas~\ref{lem:A.1}--\ref{lem:A.3}:
\(2^{(r-1)/4} = (a/r)/(b/r) = (a/r) = (-2/a)\). The Jacobi-symbol
table for odd \(a\) gives
\((-2/a) = +1 \Leftrightarrow a \equiv 1, 3 \pmod 8\).\end{proof}


\begin{remark}[Historical]
Proposition~\ref{prop:A.4} is due to Dirichlet (biquadratic
reciprocity; see Ireland--Rosen \cite[Ch.~11]{ref11} and
Lemmermeyer \cite[\S\S5.3, 5.7]{ref13}). The direct
derivation via Lemmas~\ref{lem:A.1}--\ref{lem:A.3} is included
for the reader's convenience.
\end{remark}

\begin{corollary}[Pell-form parity at Wagstaff factors]\label{cor:A.5}
At
every split Wagstaff factor
\((p, r)\) (i.e., \(r \equiv 1 \pmod 8\), \(\operatorname{ord}_r(2) =
2p\), \(p\) prime \(\geq 5\)) with Pell representation
\(r = a^2 - 2b^2\), \(a > 0\), one has \(a \equiv 1\) or \(3 \pmod 8\).
In the \(\mathbb Z[\sqrt{-2}]\) form \(r = c^2 + 2d^2\),
\(c \equiv 1\) or \(7 \pmod 8\).
\end{corollary}


\begin{proof}
By Theorem~\ref{thm:A.6} below, \(2^{(r-1)/8} \equiv \pm 1 \pmod r\),
and squaring gives \(2^{(r-1)/4} \equiv +1 \pmod r\); equivalently
2 is a fourth-power residue. By Proposition~\ref{prop:A.4}, \((-2/a) = +1\),
hence \(a \equiv 1\) or \(3 \pmod 8\). The \(\mathbb Z[\sqrt{-2}]\)
statement follows from the symmetric companion identity in the
remark above.\end{proof}


\subsection{The Octic Dichotomy, the genus formula, and the QNR Dichotomy}


\begin{theorem}[Octic Dichotomy at Wagstaff factors]\label{thm:A.6}
Let
\((p, r)\) be a split Wagstaff factor. Then \(8p \mid r-1\), and,
setting \(K'' = (r-1)/(8p)\),
\[
2^{(r-1)/8} \equiv (-1)^{K''} \pmod r.
\]
Equivalently, 2 is an eighth-power residue modulo \(r\) iff
\(v_2(r-1) \geq 4\); 2 is a fourth-power but not an eighth-power
residue iff \(v_2(r-1) = 3\). In particular, 2 is unconditionally
a fourth-power residue modulo \(r\).
\end{theorem}


\begin{proof}
By Lemma~\ref{lem:ord-2}, \(\operatorname{ord}_r(2) = 2p\), so
\(2p \mid r-1\). Combined with \(r \equiv 1 \pmod 8\) and
\(\gcd(8, 2p) = 2\), we get \(8p \mid r-1\). Hence
\((r-1)/8 = pK''\) is a positive integer, and
\[
2^{(r-1)/8} = (2^p)^{K''} \equiv (-1)^{K''} \pmod r,
\]
using \(2^p \equiv -1 \pmod r\). Since \(p\) is odd,
\(v_2(K'') = v_2(r-1) - 3\), so \(K''\) is odd iff \(v_2(r-1) = 3\),
even iff \(v_2(r-1) \geq 4\).\end{proof}


\begin{theorem}[Genus formula]\label{thm:A.7}
Let
\(r \equiv 1 \pmod 8\) be a prime and \(r = a^2 - 2b^2\) the Pell
representation with \(a, b > 0\). Then in \(\mathbb F_r\),
\[
\bigl(\tfrac{1+\sqrt 2}{r}\bigr) = (-1)^{(a-b-1)/2}.
\]
Equivalently, \(\alpha = 1+\sqrt 2\) is a square in
\(\mathbb F_r^*\) iff \(a \equiv b + 1 \pmod 4\). The right side
is invariant under the unit action
\((a,b) \mapsto (3a+4b,\, 2a+3b)\), so the formula does not depend
on the choice of representative within the Pell orbit.
\end{theorem}


\begin{proof}
Choose \(s \in \mathbb F_r\) with \(s^2 \equiv 2\); since
\(r \equiv 1 \pmod 8\) such \(s\) exists. Set \(s \equiv -a/b
\pmod r\): then \(s^2 \equiv a^2/b^2 \equiv (r + 2b^2)/b^2 \equiv 2
\pmod r\). The image of \(\alpha = 1 + \sqrt 2\) in \(\mathbb F_r\)
is \(1 + s \equiv (b - a)/b\), so
\[
\bigl(\tfrac{\alpha}{r}\bigr) = \bigl(\tfrac{(b-a)/b}{r}\bigr)
   = \bigl(\tfrac{a-b}{r}\bigr) \cdot \bigl(\tfrac{b}{r}\bigr),
\]
using \((-1/r) = +1\) and \((b/r)^2 = +1\). Set \(c = a - b\);
since \(a\) is odd and \(b\) is even, \(c\) is odd and positive
(after possibly replacing \((a,b)\) by the unit-equivalent
representative with \(a > b\)). Write \(b = 2^e b'\) with \(b'\)
odd. Then \((b/r) = (b'/r)\) (as in Lemma~\ref{lem:A.1}), and Jacobi
multiplicativity gives
\(\bigl(\tfrac{\alpha}{r}\bigr) = (cb'/r) = (r/cb')\) (the
reciprocity sign vanishes since \((r-1)/2\) is even), and
\((r/cb') = (r/c)(r/b')\). For \((r/c)\): substituting
\(a = b + c\) into \(r = a^2 - 2b^2\) yields
\(r = c^2 + 2bc - b^2 \equiv -b^2 \pmod c\), so
\((r/c) = (-1)^{(c-1)/2}\) (using \(\gcd(b,c) = 1\)).
For \((r/b')\): the Lemma~\ref{lem:A.1} reduction gives \(r \equiv a^2 \pmod{b'}\),
so \((r/b') = +1\). Combining,
\(\bigl(\tfrac{\alpha}{r}\bigr) = (-1)^{(c-1)/2} = (-1)^{(a-b-1)/2}\).\end{proof}


\begin{theorem}[QNR Dichotomy]\label{thm:A.8}
At every split Wagstaff factor
\((p, r)\), the quadratic residue character of
\(\alpha = 1+\sqrt 2\) in \(\mathbb F_r\) depends only on
\(r \bmod 16\):
\[
\bigl(\tfrac{1+\sqrt 2}{r}\bigr) =
\begin{cases}
+1 & \text{if } r \equiv 1 \pmod{16}, \text{ i.e., } v_2(r-1) \geq 4, \\
-1 & \text{if } r \equiv 9 \pmod{16}, \text{ i.e., } v_2(r-1) = 3.
\end{cases}
\]
\end{theorem}


\begin{proof}
By Corollary~\ref{cor:A.5}, \(a \equiv 1\) or \(3 \pmod 8\). Since
\(b\) is even, the joint class \((a \bmod 8,\; b \bmod 8)\) lies
in \(\{1, 3\} \times \{0, 2, 4, 6\}\). Direct computation of
\(r = a^2 - 2b^2 \pmod{16}\) on the eight cells:
\[
\begin{array}{c|cccc}
 & b\equiv 0 & b\equiv 2 & b\equiv 4 & b\equiv 6 \\\hline
a\equiv 1 & 1 & 9 & 1 & 9 \\
a\equiv 3 & 9 & 1 & 9 & 1
\end{array}
\]
So \(r \equiv 1 \pmod{16}\) iff \((a,b) \bmod 8 \in
\{(1,0),(1,4),(3,2),(3,6)\}\) and \(r \equiv 9 \pmod{16}\) iff
\((a,b) \bmod 8 \in \{(1,2),(1,6),(3,0),(3,4)\}\). By Theorem~\ref{thm:A.7},
\((-1)^{(a-b-1)/2}\) is \(+1\) on the first set and \(-1\) on the
second (check, e.g., \((1,0)\mapsto 0 \mapsto +1\) and
\((1,2)\mapsto -1 \mapsto -1\)).\end{proof}


\begin{remark}
The dichotomy is Wagstaff-specific. For arbitrary split
primes \(r \equiv 1 \pmod 8\), the QR character of \(1+\sqrt 2\)
does not have conductor 16; the field
\(\mathbb Q(\sqrt{1+\sqrt 2})\) has signature \((2,1)\) and is not
a subfield of \(\mathbb Q(\zeta_{16})\). The Wagstaff condition
collapses the higher-conductor character onto the joint cells of
\((a \bmod 8,\; b \bmod 8)\) determined by Corollary~\ref{cor:A.5}.
\end{remark}



\begin{thebibliography}{99}
\bibitem{ref1} A. O. L. Atkin and F. Morain, \emph{Elliptic curves and primality proving}, Math. Comp. \textbf{61} (1993), no. 203, 29--68.
\bibitem{ref2} Yu. Bilu, G. Hanrot, and P. M. Voutier, \emph{Existence of primitive divisors of Lucas and Lehmer numbers} (with an appendix by M. Mignotte), J. Reine Angew. Math. \textbf{539} (2001), 75--122.
\bibitem{ref3} J. Brillhart, D. H. Lehmer, and J. L. Selfridge, \emph{New primality criteria and factorizations of \(2^m \pm 1\)}, Math. Comp. \textbf{29} (1975), no. 130, 620--647.
\bibitem{ref4} J. Brillhart, D. H. Lehmer, J. L. Selfridge, B. Tuckerman, and S. S. Wagstaff, Jr., \emph{Factorizations of \(b^n \pm 1\), \(b = 2, 3, 5, 6, 7, 10, 11, 12\) Up to High Powers}, 3rd ed., Contemporary Mathematics, vol. 22, American Mathematical Society, Providence, RI, 2002, with online updates at \url{https://homes.cerias.purdue.edu/~ssw/cun/}.
\bibitem{ref5} K. S. Chua, \emph{Chebyshev polynomials and higher order Lucas--Lehmer algorithm} [preprint], arXiv:2010.02677, 2020.
\bibitem{ref6} H. Cohen and A. K. Lenstra, \emph{Implementation of a new primality test}, Math. Comp. \textbf{48} (1987), no. 177, 103--121. (Algorithm based on L. M. Adleman, C. Pomerance, and R. S. Rumely, \emph{On distinguishing prime numbers from composite numbers}, Ann. of Math. (2) \textbf{117} (1983), 173--206.)
\bibitem{ref7} R. Crandall and C. Pomerance, \emph{Prime Numbers: A Computational Perspective}, 2nd ed., Springer, New York, 2005.
\bibitem{ref8} J. Grantham, \emph{Frobenius pseudoprimes}, Math. Comp. \textbf{70} (2001), no. 234, 873--891.
\bibitem{ref9} D. R. Heath-Brown, \emph{Artin's conjecture for primitive roots}, Quart. J. Math. Oxford Ser. (2) \textbf{37} (1986), no. 145, 27--38.
\bibitem{ref10} C. Hooley, \emph{On Artin's conjecture}, J. Reine Angew. Math. \textbf{225} (1967), 209--220.
\bibitem{ref11} K. Ireland and M. Rosen, \emph{A Classical Introduction to Modern Number Theory}, 2nd ed., Graduate Texts in Mathematics, vol. 84, Springer-Verlag, New York, 1990.
\bibitem{ref12} D. H. Lehmer, \emph{On Lucas's test for the primality of Mersenne's numbers}, J. London Math. Soc. \textbf{10} (1935), no. 3, 162--165.
\bibitem{ref13} F. Lemmermeyer, \emph{Reciprocity Laws: From Euler to Eisenstein}, Springer Monographs in Mathematics, Springer-Verlag, Berlin, 2000.
\bibitem{ref14} P. Mihailescu, \emph{Primary cyclotomic units and a proof of Catalan's conjecture}, J. Reine Angew. Math. \textbf{572} (2004), 167--195.
\bibitem{ref15} P. Moree, \emph{Artin's primitive root conjecture --- a survey}, Integers \textbf{12} (2012), no. 6, 1305--1416.
\bibitem{ref16} P. Sgobba, \emph{Divisibility conditions on the order of the reductions of algebraic numbers}, Math. Comp. \textbf{92} (2023), no. 344, 2653--2688.
\bibitem{ref17} H. C. Williams, \emph{Édouard Lucas and Primality Testing}, Canadian Mathematical Society Series of Monographs and Advanced Texts, vol. 22, Wiley-Interscience, New York, 1998.
\end{thebibliography}

\section*{Supplementary material}

All computational data, verification scripts, and primality
certificates are archived on Zenodo
(\href{https://doi.org/10.5281/zenodo.19496206}{10.5281/zenodo.19496206}),
with a development mirror on GitHub at
\url{https://github.com/dolonet/wagstaff-chebyshev}.
The archive contains:
\begin{enumerate}
\item BLS primality certificates (JSON) for \(W_{2617}\),
\(W_{10501}\), and \(W_{12391}\), recording every cyclotomic
factor, BLS witness, and primality method
(Theorems~\ref{thm:W2617}--\ref{thm:W12391});
\item inert factor data for the 15{,}587{,}021 pairs \((p, r)\) with
\(G_r > 43\) (the non-trivially-closed subset of the 684{,}965{,}381
inert factors with \(r < 10^{12}\), \(p < 5 \times 10^{11}\);
\S9.5), recording \(p\), \(r\), \(G_r\),
the order parameter \(d = \operatorname{ord}_r(\omega_3)/4\),
and the Condition~(II) obstruction type;
\item danger-triple enumeration data: complete factorisations of
\(V_{114}\) and \(V_{134}\), partial factorisation of \(V_{662}\),
and the primitive-divisor survey for all 79 admissible
\(d \leq 1000\) (Theorem~\ref{thm:multi-factor});
\item scripts to reproduce the BLS proofs, the inert-factor
survey, the primitive-divisor survey (\S9.3), the
factorizations of \(U_{4d}\) and \(N_d\)
(Theorems~\ref{thm:NCT-81},~\ref{thm:pell-exclusion}), and the
danger-triple search (\S6.7); and
\item code for the Chebyshev primality test (\S9.6).
\end{enumerate}

External databases (the Cunningham project tables \cite{ref4} and
FactorDB) were used only to \emph{locate} candidate factorizations of
cyclotomic factors \(\Phi_d(2)\); every factor used in any proof was
independently re-verified for primality by APR-CL before use.
\end{document}
